\documentclass[11pt,a4paper]{article}

% --- Paquets pour la rigueur scientifique ---
\usepackage[utf8]{inputenc}
\usepackage[T1]{fontenc}
\usepackage[english]{babel}
\usepackage{amsmath, amssymb, amsfonts, amsthm}
\usepackage{geometry}
\geometry{margin=2.5cm}
\usepackage{hyperref}
\usepackage{booktabs}
\usepackage{graphicx}
\usepackage{enumitem}
\setcounter{tocdepth}{1}

% --- Métadonnées ---
\title{\textbf{The Emergence of Physical Reality within the Framework of the Projective Dynamic Logo (PDL)}}
\author{Cédric Laubscher \\ \small Independent Researcher}
\date{\today}

\begin{document}

\maketitle

\begin{abstract}

This text introduces a foundational theoretical framework, the Projective Dynamic Logo (LDP), designed to account for the emergence of physical reality from purely relational axioms, without presupposing the prior existence of space, time, or elementary particles. In this approach, reality is modelled as a network of minimal logical closures, composed of pulsating entities and their interrelations, which evolve under a principle of coherence optimisation rather than being governed by externally imposed dynamical laws.

Within these constraints, a minimal \((4,6)\) structure is derived from requirements of autonomy, reproducibility, and absence of hierarchy. This structure functions as an elementary building block for modelling the electron, whereas the proton is represented as a hierarchical architecture combining three local stationary regimes, a dynamic relational “sea”, and a finite active surface.

In the formalism adopted subsequently, the minimal \((4,6)\) structure is represented by a complete graph with four vertices whose six edges carry agreement/disagreement signs. There exists a specific assignment of signs that renders all triangular subgraphs coherent, together with a global binary pulsation \(s \leftrightarrow -s\) that preserves this coherence. This configuration provides an explicit realisation of axioms 1–3 at the elementary level and is identified with the stationary regime of an electron at rest.

This architectural framework allows several dimensionless constants to be reinterpreted as ratios of internal coherence: \(\hbar\) is associated with the minimal granularity of action; the fine-structure constant \(\alpha\) is interpreted as the fraction of the total coherence budget projected onto the surface, expressed via an active surface \(R_{\mathrm{surf}}^{(\varphi)}\) involving the mass ratio \(\mu\), the valence content \(r_{\mathrm{valence}}\), and the golden ratio \(\varphi\); \(k_{B}\) is linked to the thermal granularity associated with the dynamical fraction and a coherence leakage parameter \(\varepsilon\); and an effective gravitational coupling is construed as the metric manifestation of this minimal leakage.

Several numerical correspondences are obtained within controlled error bounds, without the introduction of additional continuous parameters. Nonetheless, some aspects of the construction remain conjectural or only partially constrained by experimental data, particularly those related to \(k_{B}\), gravitation, and cosmology. The text concludes with a speculative outline of the genesis and recycling of closures, in which Big-Bang-like episodes and extreme gravitational singularities are reinterpreted as constraint events within a pre-physical logical field. Overall, the proposal is to be regarded as a developing foundational framework rather than a fully established theory.

\end{abstract}

\newpage

% Génère la table des matières
\tableofcontents

\newpage

\section*{Introduction}
\addcontentsline{toc}{section}{Introduction}

This work is based on a distinct methodological stance: instead of postulating new physical laws on a predefined space‑time background, it adopts an upstream perspective that seeks to reconstruct physical reality from purely relational axioms. The Projective Dynamic Logo (LDP) formalises this programme by introducing a limited set of constraints on the admissible organisation of relations, such that they can generate closed and reproducible cycles. Under these constraints—logical autonomy of structures, stability of reference, absence of imposed hierarchy, and optimisation of coherence—a minimal structure, denoted \((4,6)\), emerges and assumes the role of the elementary building block of the model.

A substantial portion of the text is devoted to deriving this minimal structural configuration and examining its properties in detail, followed by the construction of more elaborate hierarchical architectures associated with nucleons, in particular the proton. The analysis then focuses on how these architectures permit several dimensionless constants—\(\hbar\), \(\alpha\), \(k_B\), and an effective gravitational coupling—to be reinterpreted as ratios characterizing internal coherence. These ratios are expressed in terms of combinatorial integers, the mass ratio \(\mu\), a leakage parameter \(\varepsilon\), and the golden ratio \(\varphi\).

This approach also yields a relational reformulation of proper time, gravitation, and selected thermal phenomena, and further develops a cosmological outline in which stable closures arise and are recurrently recycled as consequences of constraint events within a pre‑physical logical field. The work does not, however, purport to furnish a complete or self‑contained theory: some results are rigorously derived from the axioms, while others remain conjectural or only partially constrained by experimental data, and several areas—most notably detailed cosmology and out‑of‑equilibrium phenomena—are treated in a primarily exploratory fashion. The LDP should therefore be regarded as a provisional proposal for a foundational framework, intended to be extended, refined, and empirically tested through the derivation of novel, falsifiable predictions.

\newpage

\section{Limits of Conventional Physical Concepts}

Any physical theory is grounded in a set of concepts that are defined only within a specific domain of validity. When a theory is extrapolated to extreme regimes in which its fundamental concepts can no longer be coherently formulated, the difficulty no longer concerns solely experimental limitations, but the logical consistency of the concepts themselves.

In such regimes, measuring instruments may cease to operate in a meaningful way; however, what is at stake is not merely the precision of measurements, but the very possibility of attributing a well-defined meaning to the notions of space, time, energy, or physical state. The difficulty thus becomes primarily conceptual rather than empirical.

The objective is therefore not to perturbatively correct or refine existing theories at their margins, but to interrogate the conditions under which these theories can be formulated in the first place. The aim is to move upstream of any already established physical framework, toward a more primitive level of conceptualization in which the usual physical quantities are no longer presupposed.

Within such an approach, it becomes necessary to suspend the use of metric quantities and of the structural assumptions inherited from standard physical theories. Space, time, energy, and other measurable quantities can no longer be treated as primitive foundations; instead, they must be regarded as potential emergent structures, whose conditions of appearance and stability are to be elucidated. The very question of existence must then be reformulated within a framework devoid of any \emph{a priori} metric structure.

At this stage, no explicit formal hypothesis has yet been introduced. The approach remains purely logical: it consists in identifying the minimal conditions that a structure must satisfy in order to exist in a persistent manner, without invoking any prior metric postulate. The axioms that formalize these conditions will be stated explicitly in the subsequent sections.

\section{Reformulating the Question of Existence}

Once all metrics and any possibility of measurement have been suspended, the problem of existence assumes a radically different configuration: how can one distinguish “something” from nothing when no measurable property can be invoked?

In the absence of space, time, energy, or any other quantitative or qualitative parameter, no intrinsic characteristic can serve as a criterion of differentiation. Any attempt to define an entity as “that which occupies a location”, “that which endures”, or “that which possesses energy” becomes devoid of meaning, since these concepts are not yet available as operative categories. Such a situation leads necessarily to a state of indifferentiation: without a criterion, there is no basis for distinguishing one entity from another, nor even for discriminating existence from non-existence.

Within this framework, instantaneous existence does not amount to existence in the strong or robust sense. An entity that would possess no form of persistence is logically indistinguishable from nothingness: if no operation allows it to be found, identified, or recognized, it cannot be differentiated from complete absence. The fundamental problem is therefore not to specify “what exists”, but rather to determine the conditions under which something can exist persistently, that is, maintain itself as distinct from nothingness.

Under these constraints, neither a strictly immutable state nor a single, non-recurrent change can serve as a sufficient ground for existence. An immutable state, devoid of any possibility of internal or external difference, remains indistinguishable from a nothingness that is uniformly identical to itself. A single, non-repeating change, which never recurs, becomes unidentifiable: without repetition, it leaves no logical, epistemic, or ontological trace. The only coherent possibility thus resides in the emergence of a difference capable of reproducing or reiterating itself, thereby establishing the minimal structure required for persistence and, consequently, for existence in a substantive sense.

\section{Minimal alternation and informational entity}
Within this pre‑metric framework, existence cannot be grounded in substance, location, or duration, since these notions already presuppose the prior availability of spatial, temporal, or energetic metrics. What can underwrite a persistent distinction can, therefore, only be a minimal logical structure: a difference that both manifests itself and is capable of reproduction.

Given these constraints, a minimal logical alternation is posited as the simplest possible form of coherent existence. This alternation does not unfold in time or in space: in the absence of any metric, no period, frequency, trajectory, or amplitude can be ascribed to it. It has no determinable intensity; it can only be specified by the sheer fact that “something occurs rather than nothing,” and that this “something” is repeatable.

This minimal alternation can thus be characterised as an elementary informational entity. It is not “informational” in the sense of carrying semantic content, but insofar as it realises the most rudimentary form of reproducible distinction within a framework devoid of any metric. In this context, “information” does not designate a message, but the very possibility of a difference that endures.

\section{Minimal multiplicity and the constitutive role of relation}

Within a framework in which neither a metric structure nor intrinsic properties can be defined, an elementary informational entity can be distinguished from another solely by the fact of its “occurrence”. Two entities of this kind are therefore, by construction, indiscernible if their mere existence is the only available criterion of differentiation.

A strictly isolated informational entity thus remains logically undefinable. For a distinction to be operative, it is insufficient that “something happens”: a difference must be identifiable with respect to something else. Any distinction presupposes a reference, whether explicit or implicit.

In such a framework, relation is not a secondary mechanism superimposed upon already constituted entities. Rather, it is the very condition of possibility under which an informational entity can be said to exist in a definable manner. In the absence of relation, no entity can be conceptually isolated: it collapses into an undifferentiated nullity. Consequently, the existence of an informational entity can only be conceived within a minimal multiplicity, in which relation functions as the support for distinction.

A single relation, however, is insufficient to ground a stable and persistent distinction. A solitary linkage does not permit the stabilization of the reference necessary for the reproducibility of difference: any modification of one of the relata, or of the relation itself, immediately annihilates the conditions for identification. A richer, more articulated relational structure is therefore required to ensure the maintenance of distinction.

\section{Constraints on an Elementary Relational Structure}
The preceding analysis makes it possible to formulate a set of logical constraints that any admissible elementary relational structure must satisfy.

\begin{itemize}[label=-]
  \item \emph{Reproducibility of distinction}: The difference that grounds existence must be reproducible. A structure that does not allow the constitutive distinction to be instantiated repeatedly cannot support persistent existence.
  \item \emph{Stability of reference}: The referential basis with respect to which the distinction is defined must exhibit stability. A structure that depends on an external element that may disappear without leaving any trace cannot guarantee such stability.
  \item \emph{Logical autonomy}: The elementary structure must be self‑sufficient. It cannot rely on an undefined “background” to secure its internal coherence; otherwise, the distinction it implements would remain dependent on an unarticulated external context.
\end{itemize}

These constraints entail that an elementary structure must be \emph{closed}: all references required for the definition and maintenance of the distinction must be internal to the structure. An open organization, which would rely on an external frame of reference, would violate the requirement of autonomy.

Moreover, among all closed structures that meet these conditions, only those that realize them with a minimal number of relations can be retained at the fundamental level. The minimal structure is not introduced as an arbitrary postulate; rather, it is selected by the constraints themselves.

\subsection*{Fundamental Axioms of the Logical Framework}
The foregoing constraints are not technical limitations imposed a posteriori. They can be reformulated as four axioms. In what follows, a logical configuration is represented by a finite graph \(G = (V, E)\), whose vertices correspond to informational entities and whose edges encode binary agreement/disagreement relations between these entities. Each edge \(\{i,j\} \in E\) is assigned a sign \(s_{ij} \in \{+1,-1\}\), which enables the modeling of minimal alternation in the form of a discrete dynamics on these signs. Cycles of length 3 (triangles) are taken to constitute the elementary motifs of logical closure.

\subsection*{Axiom 1 – Minimal binary pulsation}
Existence is postulated to manifest as an elementary alternation between two mutually exclusive states (agreement/disagreement), which is repeatable and definable independently of any \emph{a priori} metric structure. Formally, there exists a discrete dynamical system
\[
F : \Sigma^E \rightarrow \Sigma^E, \qquad \Sigma = \{+1,-1\},
\]
such that there exists at least one sign configuration \(s^\star\) satisfying
\[
F(s^\star) = -s^\star, \qquad F(-s^\star) = s^\star.
\]
This binary pulsation does not presuppose any temporal metric; it expresses solely the existence of a logical 2-cycle in the configuration space.

\subsection*{Axiom 2 – Triangular coherence}
Any stable structure is assumed to enforce a form of logical closure: violations of this closure (open triangles, uncompensated symmetry breaking) reduce stability. For a configuration \(\sigma = (G,s)\), we consider all unordered triplets \(\{i,j,k\} \subset V\):
\begin{itemize}[label=-]
  \item the triangle \(\{i,j,k\}\) is said to be \emph{present} if the three edges \(\{i,j\}, \{j,k\}, \{i,k\}\) belong to \(E\);
  \item it is said to be \emph{coherent} if, whenever it is present, the product of the associated signs satisfies
  \[
  s_{ij}\, s_{jk}\, s_{ik} = +1
  \]
  (closed loop without uncompensated breaking);
  \item it is said to be \emph{violated} if at least one of the three edges is missing, or if the above product equals \(-1\).
\end{itemize}
We then define \(\nu(\sigma)\) as the fraction of violated triangles:
\[
\nu(\sigma) = \frac{\#\{\text{violated triangles of } G\}}{\binom{|V|}{3}},
\]
with the convention \(\nu(\sigma) = 0\) if \(|V|| < 3\). Stable structures are those that minimise \(\nu(\sigma)\).

\subsection*{Axiom 3 – Minimal completeness}
A closed structure must satisfy a completeness requirement: all references required to define and maintain the distinction must be internal to the structure itself, without recourse to any implicit external background. An elementary stationary structure is any configuration \(\sigma = (G,s)\) that satisfies:
\begin{itemize}[label=-]
  \item \(G\) is connected and complete (every pair of vertices is linked by an edge);
  \item there exists an assignment of signs \(s\) such that \(\nu(\sigma) = 0\) (all triangular subgraphs are coherent);
  \item triangular coherence forms a single block: the subgraph induced by the union of all coherent triangles is connected and contains at least two triangles sharing a common edge;
  \item the vertices of \(G\) are structurally equivalent (internal symmetry, absence of hierarchical differentiation);
  \item the dynamics \(F\) admits a binary pulsation \(s \leftrightarrow -s\) that preserves triangular coherence.
\end{itemize}
An elementary stationary structure is said to be minimal if no proper substructure satisfies these conditions. Within this framework, the first admissible minimal structure is realised for \(n = 4\), and is given by the complete graph \((4,6)\).

\subsection*{Axiom 4 – Logical optimisation}
Among all admissible configurations, only those that maximise an optimisation functional persist. This functional combines:
\begin{itemize}[label=-]
  \item minimisation of violations \(\nu(\sigma)\);
  \item maximisation of relative connectivity;
  \item realisation of closure and minimality.
\end{itemize}
We assume that the optimisation function \(\Omega(\sigma)\) depends solely on the quantities \(\nu(\sigma)\), \(\rho(\sigma)\), and an indicator of minimality, and that it satisfies:
\begin{itemize}[label=-]
  \item strict monotonic decrease with respect to \(\nu(\sigma)\) (fewer violated triangles \(\rightarrow\) greater structural stability);
  \item monotonic increase with respect to \(\rho(\sigma)\) (stronger relative connectivity \(\rightarrow\) more efficient utilisation of the internal relational budget);
  \item monotonic increase when a configuration attains the status of elementary stationary structure.
\end{itemize}
In this setting, the \((4,6)\) structure emerges as the first local maximum of \(\Omega\) at the fundamental level and thereby plays the role of elementary building block of the model.

\section{First admissible structure and introduction of logical graphs}

The first configuration not excluded by the preceding constraints corresponds to a relational structure comprising at least four relations. It remains to characterise this minimal structure in an explicit manner.

For this purpose, it is necessary to introduce a formal language capable of representing undifferentiated entities and their mutual relations without presupposing any metric or physical structure. Logical graphs provide such a language: informational entities are represented by vertices and relations by edges.

Within this framework, a closed relational structure is represented by a graph whose coherence does not depend on any vertex or edge external to the structure. The case of maximal stability, understood as a homogeneous distribution of reference, corresponds to a structure in which each entity is connected to all the others, that is, a complete graph.

The introduction of this formalism makes it possible to determine explicitly the number of entities and relations required in order to satisfy all the constraints. An admissible structure must contain at least four relations, but the number of relations cannot be considered independently of the number of entities: a closed, non-hierarchical structure requires a homogeneous distribution of relations, such that no entity occupies a privileged position.

With fewer than four entities, any attempt at such a homogeneous distribution inevitably generates degenerate or asymmetric configurations that are incompatible with the requirements of closure and logical autonomy. Starting from four entities, by contrast, it becomes possible to construct a complete graph: each entity is connected to all the others, without external dependence and without implicit hierarchy.

For a local stationary regime comprising \(n\) entities, the number of internal relations is given by
\[
r(n) = \frac{n(n-1)}{2}.
\]
For \(n = 4\), one obtains \(r(4) = 6\), which fixes the minimal structure by the ordered pair \((4,6)\).

This \((4,6)\) configuration constitutes the first case that simultaneously satisfies all the constraints and therefore defines the minimal admissible relational structure. In the formalism introduced above, this \((4,6)\) configuration is not only the first to realise a homogeneous distribution of relations; it is also the first for which there exists an assignment of signs that renders all triangles coherent \((\nu = 0)\), together with a global binary pulsation \(s \leftrightarrow -s\) that preserves this coherence, and a network of overlapping triangles that covers the entire graph. It thus provides an explicit realisation of Axioms~1 to~3 at the elementary level.

\section{Logical stationary regime and persistence}
Crossing the threshold defined by the \((4,6)\) structure does not correspond to the emergence of a novel entity in a substantial sense, nor to the introduction of an additional operative mechanism. Rather, it marks the onset of a constrained regime in which informational alternation, previously only a potential property of the system, is necessarily sustained by the relational structure itself.

More precisely, the presence of a complete \((4,6)\) graph together with a coherent sign assignment ensures the existence of a period‑2 orbit for the dynamics \(F\) on the space of signs: the agreement/disagreement alternation \(s \leftrightarrow -s\) becomes an intrinsic logical cycle of the structure. This stationary regime does not presuppose any temporal metric; it expresses solely the possibility, imposed by triangular closure, of maintaining the distinction indefinitely through an internal pulsation.

The complete closure of the minimal structure effectively imposes an internal circulation of the distinction: the difference that grounds existence can neither cease nor dissipate without compromising the global coherence of the graph. The resulting persistence is therefore not contingent but structural. It follows directly from relational saturation: as long as the \((4,6)\) structure is preserved, the distinction it supports cannot vanish without generating a contradiction.

Beyond this threshold, the distinction is no longer merely reproducible in a minimal sense; it is stabilised by the closed organisation that sustains it. This stabilisation does not depend on any external frame, pre‑given metric, or temporal or spatial reference. It is guaranteed solely by the internal coherence of the structure.

This regime can thus be characterised as a logical stationary regime: a state in which the distinction is maintained without external input, by the continuous redistribution of the relations constituting the closed structure. The internal circulation enforces an intrinsic regularity: the distinction can be maintained only in the form of a structured alternation, constrained by the closure of the whole.

This alternation manifests itself exclusively as logical periodicity, i.e., as the necessary repetition of a pattern of distinction within the closed structure, without reference to any temporal or spatial metric.

\section{Logical Optimisation Functional and Configuration Selection}

Crossing the threshold specified by the \((4,6)\) structure is not sufficient to determine a unique configuration, since a large number of more complex relational organisations remain logically admissible. The problem therefore becomes one of selection: according to which criteria can certain configurations be stably maintained, while others remain merely transient or are eliminated?

Within the framework adopted here, this selection is not ascribed to an externally imposed dynamics. Instead, it is derived from an internal optimisation principle: among all admissible configurations, only those that maximise a specific compromise between coherence, connectivity, and minimality, subject to the previously established closure constraints, are expected to persist.

To formalise this principle, we consider a configuration \(\sigma\) characterised by:
\begin{itemize}[label=-]
  \item \(n(\sigma)\): the number of informational entities involved;
  \item \(r(\sigma)\): the number of effectively active internal relations;
  \item \(\nu(\sigma)\): the number of violations of triangular coherence in the sense of Axiom~2 (open triangles, incompatible closures, uncompensated symmetry breakings);
  \item \(\rho(\sigma) = r(\sigma)/r_{\max}(n)\): the normalised relational density, where \(r_{\max}(n) = n(n-1)/2\) denotes the number of relations in a complete graph on \(n\) entities;
  \item a regime type \(\tau(\sigma)\), either stationary or dynamical, depending on whether the structure locally realises a saturated closure or only a partial connectivity.
\end{itemize}

More precisely, \(\nu(\sigma)\) is defined in terms of the triangles of the logical graph underlying \(\sigma\). For a given configuration \(\sigma\), we denote by \(G(\sigma) = (V,E)\) the graph whose vertices represent the informational entities and whose edges represent the active relations. A triplet \((i,j,k)\) of vertices forms a coherent triangle if the three edges \((i,j)\), \((j,k)\), and \((i,k)\) belong to \(E\) and if the associated pattern of agreement/disagreement alternation satisfies the closure constraint of Axiom~2 (i.e., exhibits no uncompensated “break”).

Conversely, we define a violated triangle as any triplet \((i,j,k)\) for which at least one of the three corresponding edges is absent, or for which the configuration of agreement/disagreement signs along the three edges generates an uncompensated discontinuity in the cycle. The quantity \(\nu(\sigma)\) is then defined as the total number of violated triangles in \(G(\sigma)\), normalised by the total number of possible triangles \(\binom{n(\sigma)}{3}\):
\[
\nu(\sigma)
= \frac{\#\{\text{violated triangles in } G(\sigma)\}}{\binom{n(\sigma)}{3}}.
\]

A simple illustrative case is provided by a graph with three vertices. If these three vertices are connected in a fully coherent manner (all three edges present and agreement/disagreement relations mutually compatible), then \(\nu(\sigma) = 0\). If a single edge is absent, the triplet \((i,j,k)\) forms an open triangle and one obtains \(\nu(\sigma) = 1\). In a minimal structure such as \((4,6)\), all triangles are present and satisfy the closure constraint; this entails \(\nu(\sigma) = 0\) and accounts for its maximal logical stability. The three‑vertex configuration shows that an isolated triangle can be locally coherent \((\nu(\sigma) = 0)\), but, as discussed in Section~6, it does not realise a minimal form of triangular completeness: it contains only a single cycle and does not permit the reference to be propagated through a network of overlapping triangles. The first elementary stationary structure, in the sense defined by the axioms, is therefore indeed \((4,6)\).

In what follows, we assume that the optimisation function \(F\) associated with a configuration \(\sigma\) depends exclusively on \(\nu(\sigma)\), on the relational density \(\rho(\sigma)\), and on a minimality indicator \(m(\sigma)\), according to
\[
F(\sigma)
= w_{\mathrm{coh}}\, f_{\mathrm{coh}}\bigl(\nu(\sigma)\bigr)
+ w_{\mathrm{conn}}\, f_{\mathrm{conn}}\bigl(\rho(\sigma)\bigr)
+ w_{\mathrm{min}}\, m(\sigma),
\]
where the weights \(w_{\mathrm{coh}}, w_{\mathrm{conn}}, w_{\mathrm{min}} > 0\) are fixed once and for all, and where the functional forms of \(f_{\mathrm{coh}}\) and \(f_{\mathrm{conn}}\) are chosen so as to satisfy Axioms~2 and~4. 

We then introduce an optimisation functional \(\Omega(\sigma)\) that serves as a logical selection criterion:
\[
\Omega(\sigma)
= w_{\mathrm{coh}}\!\left(1 - \frac{\nu(\sigma)}{n(\sigma)^{2}}\right)
+ w_{\mathrm{conn}}\, \rho(\sigma)
+ w_{\mathrm{min}}\, \delta_{\text{minimality}}(\sigma),
\]
where \(w_{\mathrm{coh}}, w_{\mathrm{conn}}, w_{\mathrm{min}}\) are positive weights, and \(\delta_{\text{minimality}}(\sigma)\) takes the value \(1\) if the configuration realises a closure without invoking any richer structure, and \(0\) otherwise. The first term favours configurations that minimise coherence violations \(\nu(\sigma)\) (in the sense of Axiom~2); the second term favours a degree of connectivity sufficient to guarantee the circulation of the distinction; and the third term selects organisations that remain as close as possible to the minimality prescribed by the \((4,6)\) structure.

In the present formulation of the model, we adopt for \(f_{\mathrm{coh}}\) and \(f_{\mathrm{conn}}\) normalized linear functions:
\[
f_{\mathrm{coh}}(\nu) = 1 - \nu,
\qquad
f_{\mathrm{conn}}(\rho) = \rho,
\]
so that the coherence contribution decreases linearly with the fraction of violated triangles, whereas the connectivity contribution increases linearly with the relational density.

A natural choice of weights is given by
\[
w_{\mathrm{coh}} = 1, \qquad
w_{\mathrm{conn}} = 1, \qquad
w_{\mathrm{min}} = 1,
\]
which yields
\[
\Omega(\sigma) = \bigl(1 - \nu(\sigma)\bigr) + \rho(\sigma) + \delta_{\text{minimality}}(\sigma).
\]
This choice is not intended to be unique; rather, it constitutes the simplest form compatible with Axioms~2 and~4 and is sufficient to discriminate, at a purely combinatorial level, between minimal architectures such as \((4,6)\) and composite structures that can be regarded as candidates for nucleons. A more elaborate mathematical treatment may subsequently refine this optimization functional, but the qualitative results reported here are robust with respect to the precise numerical values of the weights, provided they remain strictly positive and of comparable order of magnitude.

Logical selection can then be formulated as follows: a configuration \(\sigma\) is admissible as a persistent regime if and only if it realises a local—and, at the fundamental level, global—maximum of \(\Omega(\sigma)\) under the closure and autonomy constraints derived in the preceding sections. Configurations that fail to satisfy these constraints, or that yield a substantially lower value of \(\Omega\), remain possible as transient regimes but cannot function as stable constituents of the relational hierarchy.

When applied to elementary structures, this principle entails a particularly restrictive conclusion. For purely stationary configurations, strict minimisation of \(\nu(\sigma)\) together with the requirement of complete closure enforces a maximal relational density \(\rho(\sigma) = 1\) and precludes any internal hierarchy. The first configuration not excluded by these conditions is precisely the \((4,6)\) structure, whose relational saturation and internal symmetry maximise \(\Omega(\sigma)\) at the fundamental level. No structure with fewer entities or fewer relations simultaneously satisfies reproducibility, stability, and autonomy; furthermore, any local increase in complexity destroys this minimal character without yielding any corresponding gain in coherence.

For composite configurations, the same selection principle applies, but the trade-off between coherence and connectivity becomes more intricate. Local stationary regimes (analogous to the \((4,6)\) configuration) may coexist with dynamical regions of lower relational density, provided that the global configuration realises closure and that the associated value of \(\Omega(\sigma)\) remains maximal among the admissible architectures. Within this framework, the configurations corresponding to the proton and the neutron will subsequently be analysed: they will emerge as optima of the optimisation functional, subject to additional combinatorial constraints governing the balance between stationary and dynamical relations.

From this point onward, it becomes appropriate to investigate quantitative magnitudes associated with these optimal stationary regimes. The objective is not to impose an external metric structure, but rather to identify logical invariants that could subsequently be reinterpreted in terms of mass, electric charge, or proper frequency when these structures are instantiated in a physical context. The following section therefore develops this line of inquiry by analysing how the internal oscillation of a minimally stationary structure can be assigned a quantitative characterisation, without departing from the initially adopted pre‑metric framework.

\section{Minimal correspondence with the electron}
Within the current theoretical framework of physics, the electron emerges as the first empirical candidate compatible with the logical stationary regime associated with the \((4,6)\) structure. Observed as a strictly stable elementary entity, identical throughout the observable universe and lacking any experimentally accessible internal structure, it satisfies the stability and universality criteria derived from the relational framework.

Furthermore, in quantum field theory the electron is endowed with an intrinsic pulsation linked to its rest mass. This pulsation, commonly introduced via the Compton frequency, is independent of the electron’s trajectory and of its state of motion in any particular reference frame: it characterises the entity itself, independently of any specific geometrical description.

In this context, the proposal is not that the \((4,6)\) structure “is” an electron in a strong ontological sense, but rather that a logical stationary regime of this type admits a first coherent physical realisation in the form of the electron at rest. This identification does not rely on any additional hypothesis regarding the nature of the electron: it rests on the fact that \((4,6)\) is the first elementary stationary structure in the sense of the axioms, and that its logical pulsation \(s \leftrightarrow -s\) provides a natural candidate for the Compton alternation associated with the electron mass, via the relation \(E = \hbar \omega\). In this framework, \(\hbar\) quantifies the minimal granularity of action associated with a complete cycle of this logical pulsation. The correspondence is minimal: it consists in identifying the internal logical pulsation of the stationary regime with the physical pulsation associated with the electron mass, without invoking any supplementary mechanism.

The Compton frequency of the electron at rest can thus be interpreted as the first admissible quantised manifestation of the logical stationary regime revealed at the persistence threshold. It does not by itself constitute a proof of the logical framework, but rather its first coherent physical interpretation: an experimentally accessible quantity that can be read as the trace of an internal periodicity imposed by a closed relational structure.

Under this identification, the reduced Planck constant \(\hbar\) is no longer a purely empirical parameter; it can be interpreted as the minimal quantum of action associated with an elementary logical stationary regime. It encodes the fact that an internal pulsation such as that of the electron cannot be decomposed into arbitrarily small fractions of action: each coherence cycle carries an irreducible amount of action fixed by \(\hbar\).

More precisely, the standard relation \(E = \hbar \omega\), when applied to the proper pulsation of the electron, becomes the quantitative expression of the coherence energy required to sustain the corresponding stationary regime. In the relational framework, \(\hbar\) therefore measures the minimal “step” by which internal coherence closes onto itself, rather than a parameter characterising a pre‑existing field or particle.

\section{Energy as coherence cost}
From this standpoint, the energy associated with the pulsation of the electron should not be regarded as a fundamental property superimposed on a pre‑existing entity. Within this framework, the relation \(E = \hbar \omega\) merely renders explicit, in quantitative form, that the coherence cost associated with an elementary stationary regime is partitioned into quanta of action of magnitude \(\hbar\), each quantum corresponding to a complete cycle of the internal pulsation. This quantity can be interpreted as measuring the coherence cost required to sustain the corresponding logical stationary regime. 

Provided that the \((4,6)\) structure is preserved, the distinction it underwrites cannot be eliminated without logical inconsistency. The maintenance of this internal coherence imposes a constraint on the overall relational framework, a constraint which, when physically instantiated, appears as a well‑defined rest energy. 

At this stage, no spatial metric is presupposed, no fundamental temporality is postulated, and no supplementary internal hierarchy is introduced. Energy is defined as the quantitative expression of a particular coherence regime, rather than as a primitive quantity independent of the relational structure that supports it. 

Once the stationary pulsation is endowed with a quantitative description, energy can no longer be considered a strictly global quantity. The very possibility of multiple stationary regimes entails that the coherence cost associated with each of them can be differentiated: each elementary stationary regime must be identifiable by a characteristic value, without this identification relying on a pre‑established metric background.

\section{Possibility of Logical Localisation of Energy}
Once energy is interpreted as the coherence cost required to sustain a stationary regime, it can no longer be considered as a purely global quantity, intrinsically bound to a homogeneous background. The existence of several distinct stationary regimes implies that this cost can be differentiated from one regime to another. 

The maintenance of a particular stationary regime presupposes the imposition of a constraint associated with a specific organisation of the underlying relational structure. This constraint cannot be dissolved into an undifferentiated framework without losing its operational significance: it is necessarily anchored to the structure that realises the stationary regime. Under this perspective, energy becomes localisable in a logical sense: it is attributed to a given organisational pattern, without yet being localised in a geometrical sense. 

At this stage, localisation remains purely conceptual. It does not involve spatial position, spatial extension, or metric distance. It expresses solely the fact that coherence is not indefinitely distributable: in order for a stationary regime to be sustained, a portion of the “coherence budget” of the relational framework must be allocated to a particular structure.

\section{Coexistence of several stationary regimes}
The possibility of associating a well-defined energy with a stationary regime opens the way to the coexistence of several distinct regimes within a single relational framework. This coexistence is, however, not trivial: each regime imposes its own coherence constraints, and these constraints must remain mutually compatible if multiple regimes are to persist simultaneously. 

The coexistence of several stationary regimes thus introduces an additional requirement: each regime must be capable of preserving its internal coherence in the presence of other enduring regimes. This requirement does not yet amount to an interaction in the physical or dynamical sense; it entails neither exchange, nor force, nor any mechanism of propagation. Rather, it expresses a condition of logical compatibility: two stationary regimes can coexist in a stable manner only if the internal circulation of distinctions specific to each regime does not conflict with that of the other. 

The persistence of a stationary regime no longer depends solely on its intrinsic pulsation or on the coherence cost associated with it. It also depends on the capacity of its internal structure to be maintained without compromising the coherence of neighbouring regimes. Two regimes that are comparable in terms of pulsation and energy may nevertheless be mutually incompatible if their internal organisations do not permit stable coexistence. 

Consequently, internal stability acquires an irreducible structural character that cannot be reduced to a single scalar quantity. Compatibility or incompatibility between stationary regimes becomes a property of their internal organisational architecture, and not merely a function of their respective energy values.

\section{Towards a differential property: charge}
The coexistence of multiple stationary regimes within a single relational framework poses a specific problem: how can closed structures preserve their coherence in mutual presence without their respective internal regimes annihilating or destabilising one another? An electron, a proton, or a neutron each instantiate a stable organisation of relations; however, once any metric background that could “separate” or “isolate” these organisations is removed, their mutual compatibility is no longer conceptually guaranteed. 

Within the axiomatic framework introduced above, such compatibility is expressed in terms of triangular coherence (Axiom~2) and logical optimisation (Axiom~4). Two stationary regimes are defined as compatible if the superposition of their respective structures does not generate additional violations of triangular coherence that cannot be compensated by a global closure. Equivalently, the composition of two closed structures must remain admissible with respect to the optimisation functional: the number of violations \(\nu(\sigma)\) may increase only within bounds such that the value of \(\Omega(\sigma)\) remains maximal among the configurations accessible to the system. 

From a relational perspective, “charge” can then be interpreted as the internal signature of this differentiated form of compatibility. The logical pulsation associated with a stationary regime is not reducible to the mere fact that “an event occurs rather than none”; it exhibits an internal fine structure, determined by the distribution of agreements and disagreements over the cycles constituting the closed structure. Two regimes may thus display the same global stability while differing in the manner in which their internal alternations interlock when brought into relation.

One may define, in an abstract and purely structural manner, the charge \(q\) of a stationary regime as an internal asymmetry between the constructive and non‑constructive contributions of its relations. Relations that reinforce triangular coherence play the role of “in‑phase” contributions, whereas relations that tend to introduce uncompensated discontinuities play the role of “out‑of‑phase” contributions. In this sense, charge quantifies, in a strictly logical way, the relative excess of one of these two classes of contributions when a given stationary regime is composed with another. 

Within this framework, one may choose a reference regime whose charge is normalised to \(-1\). The minimal \((4,6)\) structure, identified in the remainder of the text with the stationary regime of the electron, constitutes a natural candidate: its complete closure and high degree of internal symmetry make it a privileged unit of comparison. A composite stationary regime whose internal pulsation exactly compensates the reference asymmetry will then realise a charge \(+1\), whereas a regime whose internal organisation is neutral with respect to this asymmetry will have zero charge. 

This definition does not treat charge as an a posteriori attribute appended to pre‑given structures. Instead, charge emerges as the consequence of two fundamental requirements: first, the necessity of preserving triangular coherence under the composition of stationary regimes (Axiom~2); second, the selection, via an optimisation principle, of configurations that realise this preservation at minimal cost in the sense of \(\Omega(\sigma)\) (Axiom~4). A strictly symmetric structure, exhibiting no exploitable internal asymmetry, would be incapable of interacting in a differentiated manner and would therefore remain neutral with respect to this property. 

In the subsequent sections, this abstract notion of charge will be instantiated in more concrete architectures. The asymmetric distribution of relations in the internal sub‑structures of a proton, together with the presence of a dynamic “relational sea”, will furnish an explicit realisation of charge \(+1\) as the signature of a minimal imbalance required for compatibility with the \((4,6)\) structure associated with the electron. Charge thus appears as the logical “cost” that a composite structure must incur in order to interact stably with an elementary structure, while preserving the global coherence of the underlying relational framework.

\section{Role of observation in the construction of the framework}

Up to this point, the development has been intentionally carried out within a purely logical framework, in which any explicit reference to physical observations has been suspended. Nevertheless, the very structure of the constraints imposed on this conceptual framework has not been devised independently of empirical input: it has been implicitly guided by robust experimental findings.

More specifically, the postulate of a minimal stationary regime reflects the existence, in established physical theory, of strictly stable entities such as the electron, whose intrinsic properties remain invariant across all experimentally accessible observational contexts. Similarly, the requirement to distinguish stationary regimes that coexist without mutual amalgamation mirrors the empirical evidence that multiple particles can be present simultaneously without ceasing to be individually identifiable.

Observation, therefore, does not operate as a direct generator of axioms, but rather as a retrospective selection filter: among the set of logically conceivable frameworks, only those are retained that permit at least a minimal correspondence with empirically observed structures. It is precisely because experimental physics discloses entities that are stable, localisable, and mutually distinguishable that the logical construction is constrained to give rise to stationary regimes, localisable energy distributions, and well-defined charges.

In this sense, the epistemic role of observation is twofold. First, it enables the entire undertaking by revealing the persistent features that the logical framework must be capable of reconstructing. Second, it furnishes the selection principle among competing logical constructions: only those frameworks that admit a coherent and systematic alignment with the empirically observed world can be regarded as viable candidates for a description of the emergence of physical reality.

\section{Necessity of a pre‑metric separation}

The introduction of charge as an internal differential property enables the characterization of the mutual compatibility or incompatibility of stationary regimes coexisting within a common relational framework. However, as long as this framework remains entirely undifferentiated, there is, at the purely logical level, no constraint preventing mutually incompatible regimes from overlapping without any discernible distinction.

Such a situation would contradict the requirement of persistence. If two incompatible stationary regimes were allowed to merge, the internal coherence of each regime would be disrupted, and no stable structure could be sustained. The coexistence of stationary regimes bearing distinct charges therefore necessitates the existence of a minimal separation between incompatible configurations.

This separation cannot yet be interpreted as a spatial distance, since no metric structure has been specified. Nor can it be described as the consequence of an interaction or a dynamical process, as neither forces nor propagation mechanisms have been defined at this stage. It constitutes a pre‑metric separation, determined exclusively by coherence constraints: certain combinations of regimes are ruled out, not because they “repel” one another, but because they cannot occupy the same relational organisation without engendering logical inconsistency.

\section{Towards a Minimal Metric}

Once multiple stationary regimes coexist—distinguished by their charges and separated by coherence constraints—it becomes necessary to introduce a formal apparatus capable of representing these separations without modifying their logical status. 

Such an apparatus does not generate the separations; rather, it yields a systematic reading of them. It does not constitute an additional physical structure, but a minimal instrument designed to render comparable relations that are already determined by compatibility constraints. The objective is to introduce a minimal metric which does not, at this stage, organise a continuous geometric space, but which nonetheless makes it possible to compare and order the separations between stationary regimes. 

This minimal metric introduces neither geometry in the classical sense, nor continuity, nor dynamics. It serves solely to distinguish situations in which two stationary regimes may share the same logical neighbourhood from those in which they must be kept distinct. It furnishes a language for asserting that one configuration is “more constrained” than another, without yet assigning spatial distances or temporal durations in the physical sense. 

It is crucial to stress that this minimal metric has no fundamental ontological status. It is not posited as a pre‑existing structure of reality, but is introduced as a constrained response to a representational problem: in its absence, it would be impossible to describe in an ordered and comparative manner a collection of stationary regimes differentiated by their charges and by their compatibility relations.

\section{Composite structures and minimal hierarchy}
The minimal metric furnishes a representational framework for multiple coexisting stationary regimes, distinguished by their respective charges and separated by coherence constraints. Empirical observation, however, reveals the existence of entities whose rest mass substantially exceeds that of the electron while exhibiting comparable stability. These entities cannot be reduced to simple aggregates of independent electrons: they display an internally organised structure, characterised by a coherence that surpasses the mere additive contribution of their constituents. 

Such global coherence cannot be accounted for by a naive juxtaposition of elementary stationary regimes, each preserving its own complete closure. If every sub‑structure were to maintain a maximal closure of type \((4,6)\), the resulting total relational density would reach a level incompatible with the presence of any internal hierarchy; the composite structure would then be rendered indecomposable. 

The existence of a stable composite organisation therefore requires an architecture in which relational saturation is non‑uniform: certain regions attain maximal closure, whereas others preserve only partial connectivity. This heterogeneous distribution of relational density defines the minimal hierarchy of a composite structure. 

In a minimal structure such as the \((4,6)\) configuration, closure is complete: all relations necessary for coherence are internal, and no coherence circulation remains unsupported. Such a structure exhibits no logical leakage and realises a maximal saturation consistent with stability. 

A hierarchical composite structure is subject to different constraints. To prevent global saturation, it must sustain some regions in a local stationary regime and others in a dynamical regime. This internal asymmetry precludes, by construction, perfect closure at the global scale: a portion of the coherence circulation cannot be fully confined within the closed sub‑structures. It must instead traverse the collective relational environment. 

One may then introduce the notion of coherence leakage: a minimal fraction of the circulation required to maintain the structure fails to close entirely within the internal hierarchy and propagates through the surrounding relational framework. This leakage does not correspond to a loss in the energetic sense; rather, it expresses the logical impossibility, for a hierarchical composite structure, of simultaneously realising both maximal local closure and maximal global closure.

\section{The proton as a minimal composite structure: a two‑level architecture}
The proton cannot be represented as a simple juxtaposition of three identical instances of the minimal \((4,6)\) structure. Instead, it realises a two‑level hierarchical architecture, consisting of local stationary regimes—naturally identifiable with valence quarks—and a dynamically fluctuating relational “sea” that interconnects these local cores. 

The stationary substructures function as local centres of coherence: each of them maximises its internal closure, analogously to the \((4,6)\) configuration, but at a higher level of compositional complexity. In contrast, the sea constitutes a region of reduced relational saturation, in which the relations are not fully closed but are sufficient to sustain and transmit coherence between the quark cores. 

This structural asymmetry is dictated by the axioms of coherence and optimisation. If all substructures were simultaneously saturated in the same manner as the \((4,6)\) unit, the resulting global relational density would inhibit any effective propagation of coherence among them. A minimal hierarchy is therefore required: local maximisation of closure at the level of the cores, combined with only partial closure at the global scale, with the sea acting as the mediating domain that enables coherence transfer within the composite system.

\subsection{Combinatorial constraints on valence quarks}
Local stationary regimes must satisfy the minimal completeness condition imposed by Axiom~3: the number of elementary entities constituting a quark must be a multiple of 4, so that its internal organisation can be represented as a superposition of logical tetrads. Consequently, for a generic quark one has:
\[
n_{\mathrm{quark}} \in \{4, 8, 12, 16, 20, 24, 28, \ldots\}.
\]
In the case of a closed stationary configuration, each quark is modelled by a complete graph \(K_n\), with:
\[
r_{\mathrm{quark}} = \frac{n_{\mathrm{quark}}\bigl(n_{\mathrm{quark}} - 1\bigr)}{2},
\]
where \(r_{\mathrm{quark}}\) denotes the number of internal relations (edges) between the entities (vertices) of the quark.

The internal completeness requirement is further constrained by a minimal asymmetry condition: for a non-vanishing net charge to arise, the three valence quarks cannot possess strictly identical internal architectures. Two of them must share the same combinatorial structure (identified with the up quarks), while the third (identified with the down quark) must realise a slightly distinct internal organisation.

This controlled breaking of internal symmetry constitutes the necessary condition for the global configuration to maximise the optimisation functional \(\Omega(\sigma)\) while generating a total charge \(+1\), in a manner compatible with the charge \(-1\) associated with the \((4,6)\) structure modelling the electron.

\subsection{Rationale for the choice \(n_u = 24\) and \(n_d = 28\)}
Under these constraints, we examine a discrete set of candidate configurations in which the valence quark parameters obey
\[
n_u, n_d \in \{4, 8, 12, 16, 20, 24, 28, \ldots\},
\]
with two identical up quarks and one distinct down quark. For each admissible pair \((n_u, n_d)\), the associated composite configuration is assessed via the optimisation functional \(\Omega(\sigma)\) defined in Section~10, which incorporates the number of triangular coherence violations \(\nu(\sigma)\), the global relational density, and a criterion of hierarchical minimality.

As representative examples, the choices \((n_u, n_d) = (12,16)\) and \((16,20)\) do yield valence quark states constructed as superpositions of logical tetrads (3 and 4 tetrads, respectively). However, these configurations are ruled out for several reasons:
\begin{itemize}[label=-]
  \item the global relational density remains too close to saturation, thereby inhibiting the emergence of a sufficiently dynamic relational sea: the value of \(\rho(\sigma)\) is excessively high, and the contribution \(f_{\mathrm{conn}}(\rho)\) fails to offset the concomitant increase in \(\nu(\sigma)\);
  \item the resulting distribution of effective charges does not permit the realisation of a \(u,u,d\) configuration with total charge \(+1\) without introducing marked asymmetries, which in turn drives \(\nu(\sigma)\) beyond the threshold compatible with structural stability.
\end{itemize}

Conversely, larger pairs such as \((n_u, n_d) = (28,32)\) and beyond provide sufficient capacity to accommodate a substantial relational sea, but only at the expense of a total number of relations \(r_{\mathrm{total}}\) that significantly exceeds \(1.1 \times 10^{4}\). Such architectures are excluded by the minimality criterion, for which \(\delta_{\text{minimality}}(\sigma) = 0\), and the corresponding value of \(\Omega(\sigma)\) is markedly lower than that obtained for the configuration \((24,28)\).

Within this framework, the pair \((n_u, n_d) = (24,28)\) thus appears as the first combination that maximises \(\Omega(\sigma)\) under the imposed combinatorial constraints: divisibility of \(n_u\) and \(n_d\) by 4, minimal asymmetry required to generate a net charge \(+1\), presence of a non-negligible relational sea, and a stationary fraction \(\sim 9/10\) compatible with the empirically inferred share of valence energy in the proton mass.

In the course of this combinatorial exploration, the pair \((n_u = 24, n_d = 28)\) emerges as the first configuration that simultaneously satisfies: (i) internal completeness of each quark, expressed by the requirement that \(n_u\) and \(n_d\) be multiples of 4, (ii) minimal asymmetry sufficient to yield a net charge \(+1\), and (iii) a maximal value of \(\Omega(\sigma)\) among the configurations examined. The values
\[
n_u = 24 = 6 \times 4, \qquad n_d = 28 = 7 \times 4
\]
may therefore be interpreted as the outcome of a constrained selection process rather than as arbitrary parameters: they correspond to superpositions of six and seven logical tetrads, respectively, which optimise the internal coherence of the proton within this relational framework.

On this basis, the numbers of stationary relations associated with the up and down quarks are
\[
r_u = \frac{24 \times 23}{2} = 276, \qquad
r_d = \frac{28 \times 27}{2} = 378,
\]
thereby setting the stage for the subsequent detailed analysis of the valence contributions and of the relational sea presented in the following subsections.

\subsection{Quantification of the relational sea}
The quark sea is modeled as a sparse graph with minimal connectivity, characterized by the following counting rule:
\[
R_{\mathrm{sea}} = 2\, n_{\mathrm{sea}}.
\]
For the proton, application of this rule yields:
\[
R_{\mathrm{sea}} \simeq 10\,087.
\]
This term quantifies the dynamical (non‑stationary) component of the total relational coherence.

\subsection{Total relational content and stationary/dynamic fractions}

The total number of internal relations characterising the relational architecture of the proton is
\[
r_{\mathrm{total}} = r_{\mathrm{valence}} + R_{\mathrm{sea}}
= 930 + 10\,087
= 11\,017.
\]
From this, one immediately obtains the stationary and dynamical fractions:
\[
f_{\mathrm{stat}} = \frac{930}{11\,017} \simeq 9\%, \qquad
f_{\mathrm{dyn}}  = \frac{10\,087}{11\,017} \simeq 91\%.
\]
These values are in qualitative agreement with the empirically established result that approximately \(91\%\) of the proton’s energy arises from internal dynamical degrees of freedom (gluon field and quark sea), rather than from the bare rest masses of the valence quarks. 

Within the LDP framework, the approximate \(9/91\) partition emerges from a hierarchical relational architecture determined by Axioms~2–4 (minimisation of violations, maximisation of connectivity, and enforcement of a minimality constraint) together with an optimisation functional whose weights are chosen to encode the predominance of dynamical regimes. It does not result from a direct fit to a specific experimental value. 

The close correspondence with QCD-based estimates can therefore be interpreted as an a posteriori validation of the LDP‑derived architecture, instead of being the by‑product of a parametric adjustment tuned to empirical data. 

The integers introduced here are not treated as free parameters; rather, they are selected under a dual set of constraints:
\begin{itemize}[label=-]
  \item consistency with the axioms (multiplicity by 4, realisation as complete graphs or minimal sparse graphs, and minimisation of \(\nu\));
  \item reproduction of dimensionless ratios known from QCD (in particular the \(\approx 9/91\) split, and the relevant orders of magnitude for radii and reduced masses).
\end{itemize}
Numerical explorations of alternative integer combinations indicate that these either require substantially larger relational sizes or entail an accumulation of violations incompatible with Axiom~4. This behaviour supports the claim that the adopted architecture is not arbitrary, but instead singled out by the combined constraints. 

In the graphical formalism introduced above, the underlying structures are represented as large signed graphs in which three local stationary closures play the role of effective valence quarks. The integers \(930\), \(10\,087\), and \(11\,017\) enumerate, respectively, the stationary relations associated with these closures, the dynamical relations constituting the “relational sea”, and the total number of internal relations. Their specific values are obtained through a joint optimisation of triangular coherence (minimal \(\nu\)), relational density \(\rho\), and hierarchical minimality, all subject to the constraints encoded in Axioms~2–4.

\subsection{Coherence leakage}
The hierarchical structure cannot simultaneously achieve perfect (maximal) closure of all its sub‑systems and perfect global closure of the entire system. Consequently, a minimal coherence leakage necessarily arises: a very small fraction of the relational circulation required to ensure the stability of the proton does not close entirely within its \(11\,017\) internal relations, but instead propagates into the surrounding relational framework. 

We therefore introduce a leakage parameter \(\varepsilon\), which quantifies the relative deviation between an ideal, fully closed hierarchy and the effective degree of closure actually realised by the protonic structure. For a minimal composite structure such as the proton, this parameter is expected to be very small. It represents a minimal logical violation of perfect closure, sufficient to establish a non‑trivial coupling between the internal hierarchy and the external relational environment. 

In the formal language of the axioms, \(\varepsilon\) can be interpreted as the residual fraction of coherence that cannot be captured by a strict minimisation of \(\nu(\sigma)\) restricted to the graphs internal to the proton. It thus measures the discrepancy between (i) the optimum that would be attained under an ideal global closure, in which all triangular configurations across all hierarchical levels are coherent, and (ii) the optimum that is actually attainable under the minimality and hierarchical constraints encoded in \(\Omega(\sigma)\). 

In the subsequent analysis, this parameter will play a central role: it will allow us to establish a quantitative relation between the internal structure of the proton and the curvature properties of the emergent metric, and thereby to identify an effective gravitational coupling constant.

\section{Active surface of the proton and coupling with the electron}
The relational description of the proton distinguishes two principal contributions: \(930\) stationary relations associated with the three valence quarks, and a dynamical quark sea comprising approximately \(10\,087\) relations. However, not all of these relations play an equivalent role with respect to an electron: during the coupling process, only a subset of the protonic structure effectively participates in the interaction.

To characterise this subset, we introduce the notion of the \emph{active surface} of the proton. By active surface we mean the minimal envelope of relations that actually mediate the logical coupling with an electron described by the \((4,6)\) structure. This surface is distinct both from the full set of internal protonic relations and from the simple sum of valence-quark relations. Rather, it is defined as the subset of relations that (i) bear the \(+1\) effective charge signature and (ii) remain accessible to the stationary regime of the electron.

Within the formalism of signed graphs, this coupling is implemented via mixed triangles that connect vertices of the protonic structure with vertices of the \((4,6)\) structure associated with the electron. The relations that constitute the active surface are precisely those that participate in \emph{coherent} mixed triangles, i.e. triangles which, when combined with the internal sign structure of \((4,6)\), do not introduce additional violations of triangular coherence. The effective proton charge \(+1\) with respect to the electron is thereby represented as a global asymmetry favouring such constructive contributions within the ensemble of accessible mixed triangles.

To estimate this active surface, it is convenient to separate the contributions of the two up quarks (\(u\)) and the down quark (\(d\)). The up quarks, each carrying charge \(+2/3\), contribute to surface coherence through their \(276\) stationary relations, whereas the down quark, with charge \(-1/3\), contributes via its \(378\) relations. A natural estimate of the “surface charge content” is then obtained by weighting these relational contributions by the corresponding quark charges, which leads to an effective relational combination of the form
\[
R_{\mathrm{raw}} \sim 2\, r_u + r_d
= 2 \cdot 276 + 378
= 930,
\]
in numerical agreement with the total number of valence-quark relations.

Within the framework of closure theory, the internal \(u\) and \(d\) blocks do not fulfil the same structural function as the \((4,6)\) elementary brick. They do not constitute autonomous elementary stationary structures, but rather local stationary closures: densely interconnected subgraphs—of cardinality \(24\) for \(u\) and \(28\) for \(d\)—in which triangular coherence is strongly maximised and an internal pulsation, consistent with the axiomatic constraints, can be sustained. Their complete coherence can nevertheless only be guaranteed when they are embedded in the full hierarchical architecture of the proton, which encompasses the relational sea and the global closure conditions.

The internal \(u\) (24 entities, 276 stationary relations) and \(d\) (28 entities, 378 stationary relations) blocks are thus interpreted as local stationary closures within the proton: highly connected regions in which triangular coherence is significantly maximised and an internal pulsation is realisable, but which, taken in isolation, do not satisfy all the minimal completeness criteria defined for the \((4,6)\) brick. Their logical stability is therefore dependent on the global hierarchical organisation of the proton (valence sector, sea, surface, and leakage).

Moreover, not all valence relations are equally effective in mediating interactions with the exterior. A subset of these relations is devoted to maintaining the internal coherence of the quark blocks and does not couple directly to the interface with the electron. By incorporating this internal hierarchy (locally closed quarks, quark sea, and global closure constraints), one obtains a reduced effective surface comprising on the order of
\[
R_{\mathrm{surface}} \simeq 500
\]
active relations. This effective quantity represents, in combinatorial terms, the size of the subset of relations that actually participates in coherent mixed triangles with the \((4,6)\) structure. It is smaller than the bare valence content \((930)\), because a fraction of the stationary relations within the \(u\) and \(d\) blocks remains committed to preserving their internal coherence and does not project directly onto the logical interface with the electron.

The distinction between total relations and surface relations is crucial for the subsequent analysis. The fine‑structure constant \(\alpha\) can be reinterpreted as quantifying the relative efficiency of this coupling, through a comparison of the active proton surface, the stationary regime of the electron, and the mass ratio \(\mu = m_p/m_e\). The next section develops this perspective to propose a relational reinterpretation of \(\alpha\), and to show how the order of magnitude \(\alpha^{-1} \simeq 137\) may emerge from the internal architecture of the proton together with the minimal character of the electron.

\section{Proton Compton wavelength and closure of the shape}

The relational structure of the proton has so far been characterised without invoking any \emph{a priori} spatial metric. Experimentally, however, the proton is observed as an entity endowed with an effective radius and a Compton wavelength \(\lambda_p\), associated with its rest mass. Within the present framework, this wavelength is not regarded as an independent empirical input, but as the first metrical translation of a global stationary regime associated with the closed shape.

Let \(R_p\) denote an effective logical radius of the proton. Elastic scattering measurements typically yield \(R_p \simeq 0.84 \times 10^{-15}\,\mathrm{m}\). To first approximation, a stationary wave propagating along the proton surface can be modelled as following a great‑circle trajectory. Assuming that a single wave achieves closure of the shape over a half‑turn, the corresponding arc length is \(\pi R_p\). The recommended Compton wavelength of the proton is \(\lambda_p \simeq 1.32 \times 10^{-15}\,\mathrm{m}\).

The ratio \(\lambda_p / (\pi R_p)\) is then close to \(0.5\), suggesting that, instead of a single closing wave, two counter‑circulating waves propagate on the logical surface of the proton, each performing a half‑turn per period. The choice of a great‑circle propagation should be interpreted as a minimal model for a surface mode among a broader family of geometrically admissible configurations. It is not intended to exhaust all solutions compatible with a closed shape, but rather to provide a simple scenario for assessing the consistency between logical closure and its first metrical translation:
\[
2\pi R_p \simeq 2 \times 1.32 \times 10^{-15}\,\mathrm{m},
\]
which, for \(R_p \simeq 0.84 \times 10^{-15}\,\mathrm{m}\), yields numerical agreement at the level of better than one percent. The relative deviation between the radius inferred from \(\lambda_p\) via the relation \(\lambda_p \simeq 2\pi R_p\) and the radius extracted from scattering data may be interpreted as a slight “breathing” of the shape, associated with the internal dynamics of the quark–gluon sea.

In this perspective, the “Compton wavelength” does not introduce a novel fundamental quantity: it merely reuses the standard relation \(\lambda = h/(mc)\) (or \(\hbar/(mc)\), depending on convention), applied to the proton rest mass, and functions as a point of contact between the internal frequency and an effective metric scale. The numerical agreement at the sub‑percent level must therefore be interpreted as a plausibility check of the proposed surface‑wave model, rather than as a unique or definitive derivation of proton geometry.

This numerical coherence does not constitute a proof, but rather an indication that the proton Compton wavelength can be interpreted as the geometric manifestation of a stationary surface mode imposed by the global stationary regime. In this view, space ceases to be a neutral, pre‑given container and instead appears as the manner in which a stable logical closure projects itself into a metrical description.

\section{Proper time as cycle counting}

Interpreting the Compton wavelength of the proton as the manifestation of a stationary surface wave makes it possible to assign an analogous status to the Compton frequency associated with its rest mass. This frequency does not describe a motion in an already given time; it expresses the number of pulsation cycles required to maintain the global stationary regime. 

\subsection*{Proper time and relational density}
Within the relational framework, the proper time associated with a given structure is not conceived as an external coordinate, but rather as a measure of the number of coherence cycles that the structure must execute in order to sustain its existence. To connect this notion with a more geometrical representation, it is useful to introduce the relational density \(\rho\), defined as the number of active relations per unit of logical volume.

One may then express, in schematic form, that the proper‑time element \(\mathrm{d}\tau\) associated with a region of relational density \(\rho\) is related to a reference pulsation increment \(\mathrm{d}t_{0}\) by a relation of the type
\[
\mathrm{d}\tau \sim \sqrt{\rho}\,\mathrm{d}t_{0}.
\]
In a region of high relational density (e.g.\ dense matter, strong concentration of hierarchical structures), the effective value of \(\mathrm{d}\tau\) per unit of \(\mathrm{d}t_{0}\) is reduced: the structure must allocate a larger fraction of its cycles to the maintenance of its internal coherence. For an external observer, this manifests itself as a dilation of proper time. Conversely, in a region of low relational density (approaching a logical vacuum), the structure expends fewer cycles on preserving its coherence, and proper time elapses more rapidly with respect to the same reference increment \(\mathrm{d}t_{0}\).

This qualitative dependence reproduces the general pattern of empirically established time‑dilation phenomena: proper time is slowed in regions of high matter or energy density and accelerated in “emptier” regions. The departure from general relativity lies in the interpretation of the curvature of the emergent metric: here, it is regarded as the effective trace of the distribution of relational density \(\rho\) and of the corresponding coherence cost, rather than as an autonomous fundamental geometrical structure.

In the case of the electron, the Compton frequency \(\nu_e\) provides an intrinsically very high intrinsic cadence: each period corresponds to a complete re‑affirmation of the internal coherence of the minimal \((4,6)\) structure. The proper time associated with the electron can then be defined as the counting of these cycles, i.e.\ as the number of iterations through which the structure confirms its stability.

This definition does not alter the standard relativistic notion of proper time; instead, it recasts it in terms of internal cadence by using the relation \(E = \hbar \omega\) applied to the rest mass. In the frame in which the electron is at rest, the “cycle counting’’ associated with its Compton pulsation thus coincides with the conventional measurement of its proper time, thereby ensuring compatibility with the usual relativistic description.

For a composite nucleon such as the proton, the situation is more intricate. The global frequency associated with its rest mass defines an average cadence, but this cadence emerges from the hierarchical coupling between:
\begin{itemize}[label=-]
  \item the pulsations of local stationary regimes (valence quarks);
  \item the more weakly constrained dynamics of the collective regime (quark sea);
  \item and the global stationary regime that determines the Compton surface wave.
\end{itemize}

Within this framework, the fine‑structure constant \(\alpha\) can be interpreted as a ratio between the total internal coherence of the proton and the portion of this coherence that remains available at its effective surface for external couplings. The logical surface on which the global Compton wave manifests is not a mere geometric boundary; rather, it is a locus that concentrates a finite number of active relational degrees of freedom, originating from the internal hierarchy composed of valence quarks and the surrounding quark sea, which remain accessible for coupling with external structures such as the electron. 

The dimensionless value \(\alpha \approx 1/137\) thus expresses the fraction of the proton’s “coherence budget” that is projected into these surface relations, relative to the internal saturation required to ensure its hierarchical closure. In this relational interpretation, \(\alpha\) quantifies the “fineness” of electromagnetic coupling not as a constant imposed externally, but as the quantitative signature of the share of surface coherence that a proton can allocate to its interactions without compromising its internal stability. 

The proper time of a proton can accordingly be conceived as the counting of the cycles of this global pulsation, each cycle integrating the contributions of all sub‑structures that contribute to hierarchical coherence. Two structures with identical rest masses may therefore exhibit the same global frequency, while differing in the density of internal pulsations they mobilize to maintain their respective coherence. 

In this perspective, the metric time of established physics appears as an effective average over these hierarchical countings, rather than as an independent fundamental parameter. The dilation of durations observed in relativity can then be reinterpreted as a modification of the coherence conditions imposed on these stationary regimes when they are subjected to extreme relational configurations, rather than as a deformation of an abstract temporal substrate. This constitutes a conceptual re‑interpretation and does not contest the quantitative predictions of special and general relativity: time‑dilation effects continue to be described by the same mathematical relations, but are here understood as variations in the coherence budget available to stationary regimes, rather than as anomalies of a pre‑given fundamental time.

\section{The fine-structure constant as an indicator of surface coherence}

In the standard formulation of quantum electrodynamics, the fine-structure constant \(\alpha \simeq 1/137\) is introduced as a dimensionless parameter that quantifies the strength of the electromagnetic interaction. Within the relational framework developed here, it can instead be interpreted as a quantitative indicator of a coherence ratio between the internal organisational structure of the proton and the fraction of this structure that remains accessible at its effective surface for external couplings.

In a minimal \((4,6)\)-type structure, such as the one attributed to the electron, no analogous constant is meaningfully defined: closure is complete, there is neither an internal hierarchy nor a logically distinct surface separable from the overall coherence volume. Only upon crossing a hierarchical threshold—marked by the emergence of a composite structure such as the proton, which combines local stationary regimes, a relational “sea,” and coherence leakage—does a relational “skin” become distinguishable from a saturated core, thereby rendering an internal-to-surface coherence ratio a well-defined concept. From this standpoint, \(\alpha\) can be regarded as a secondary manifestation of this change of organisational regime: it encodes the proportion of coherence that the composite structure can maintain at its surface while preserving internal stability.

\subsection{Internal coherence and logical surface of the proton}

The proton, previously described as a minimal composite structure, can be regarded as comprising three local stationary regimes of quark type together with a collective relational sea that guarantees global cohesion without reaching a saturation threshold. The total number of internal relations, obtained via a detailed relational enumeration, is of the order of
\[
r_{\mathrm{total}} \simeq 1.1 \cdot 10^{4},
\]
which quantifies the global coherence cost required to sustain the proton’s hierarchical organization.

However, not all of these relations are equivalent from the standpoint of external couplings. A fraction of them remains confined within the local closures associated with the valence quarks and the quark sea, whereas a subset constitutes an active logical surface: a finite ensemble of relations through which the proton can establish stable couplings with external structures, in particular with a stationary regime of electron type.

This subset of surface relations does not correspond to a geometric boundary in the classical sense, but rather to a relational “skin” that remains available once internal coherence has been saturated. Combinatorial analyses associated with the internal structure of the proton yield an effective number of surface relations on the order of
\[
R_{\mathrm{surf}} \simeq 500,
\]
which will be adopted in the following as the characteristic measure of the active surface.

\subsection{Mass ratio and coherence scale}
The ratio between the mass of the proton and that of the electron,
\[
\mu = \frac{m_p}{m_e} \simeq 1836,
\]
represents, within the relational framework, the relative difference in coherence cost between the composite stationary regime (proton) and the minimal elementary stationary regime (electron). It is not a purely empirical parameter: this ratio encodes the manner in which two distinct forms of hierarchical closure, characterized by markedly different internal complexities, are mapped onto the quantitative language of rest mass. 

The electron is associated with the minimal \((4,6)\) structure, whose exact closure determines an intrinsic pulsation that is directly expressed through its Compton frequency. By contrast, the proton realises a composite hierarchical closure comprising a stationary component of the order of a few percent and a dynamic component of approximately \(91\%\), sustained by the collective relational background (or “relational sea”). The ratio \(\mu\) thus quantifies the disparity between these two coherence regimes, once they are expressed in the common units defined by \(\hbar\) and \(c\).

\subsection{Relational definition of \(\alpha\)}

In the previous section, the active surface of the proton was estimated to be \(R_{\mathrm{surf}} \simeq 500\) relations, representing the minimal envelope through which the proton can establish a stable coupling with an electron. This estimate now makes it possible to reinterpret the fine-structure constant as a coherence ratio between this active surface and the global internal architecture of the proton.

The objective here is not to derive a unique theoretical prediction for \(\alpha\), but rather to show that a single protonic architecture—already constrained by QCD and by the axioms of LDP—naturally yields a value close to the empirical \(\alpha\), without recourse to continuous parameter tuning.

Within this framework, the fine-structure constant can be expressed as the dimensionless ratio
\[
\alpha^{-1} \simeq \frac{R_{\mathrm{surf}}^{2}}{\mu},
\]
where \(R_{\mathrm{surf}}\) denotes the number of active surface relations of the proton, and \(\mu\) is the proton–electron mass ratio.

For \(R_{\mathrm{surf}} \simeq 500\) and \(\mu \simeq 1836.15\), one obtains
\[
\alpha^{-1}
\simeq \frac{500^{2}}{1836.15}
\approx 136.15,
\]
which agrees, to better than one percent, with the experimental value \(\alpha^{-1} \simeq 137.036\).

This result may be compared with the experimental value \(\alpha_{\mathrm{exp}} \simeq 1/137.036\). The relative discrepancy, of order one percent, is compatible with the approximations involved both in the counting of surface relations and in the internal hierarchical modeling of the proton. It suggests that interpreting \(\alpha\) as a surface-coherence ratio is, at the very least, quantitatively plausible.

The above relation does not constitute an independent prediction of \(\alpha\), as it explicitly relies on the empirical value of the mass ratio \(\mu\). Rather, it offers a reinterpretation of the fine-structure constant in terms of surface coherence, numerically consistent with the observed value, without yet establishing that this value is uniquely determined within the proposed axiomatic framework.

The proximity of the obtained value for \(\alpha^{-1}\) (relative error \(\sim 10^{-3}\)) does not arise from any continuous fine-tuning, but from an integer combination fixed by the relational architecture (\(\approx 500\) active relations out of 930) together with the mass ratio \(m_p/m_e\), with no freedom associated with additional adjustable parameters.

\subsection{Interpretation: coupling fineness and hydrogen stability}
The relational interpretation of \(\alpha\) does not consist in introducing an additional fundamental constant, but rather in reinterpreting an already established constant of physics as a quantitative measure of the fraction of surface coherence. A significantly larger value of \(\alpha\) would imply that a greater portion of the proton’s internal coherence budget is made available at its surface; the resulting enhancement of the coupling to the electron would then be excessive, potentially jeopardising the stability of bound states such as the hydrogen atom. Conversely, a substantially smaller value of \(\alpha\) would correspond to a logical surface whose activity is too weak to support the formation of stable bound states.

The empirically observed value of \(\alpha\) thus emerges as the unique compromise between two competing constraints:
\begin{itemize}[label=-]
  \item preserving the internal hierarchical closure of the proton by allocating the majority of coherence resources to the integrity of its composite structure;
  \item maintaining a sufficiently active relational “skin” to enable the stabilisation of an electron-type stationary regime at a finite distance.
\end{itemize}
This interpretation is compatible with the fact that the binding energy of the hydrogen atom, derived within the standard Coulomb framework, can be reformulated in terms of an effective potential in which \(\alpha\) enters as a factor representing the fraction of surface coherence available for coupling, without altering the empirical value \(-13.6\,\mathrm{eV}\) at the Bohr radius.

\subsection{\texorpdfstring{\(\alpha\), \(\hbar\)}{alpha, hbar} and the granularity of couplings}

Within the relational framework, the reduced Planck constant \(\hbar\) is interpreted as the minimal quantum of action associated with an elementary logical stationary regime. Each cycle of internal pulsation carries a discrete unit of action of magnitude \(\hbar\), and the proper energy of a structure is then expressed as \(E = \hbar \omega\), provided that a physical realisation of this internal pulsation is available.

In the same spirit, the fine-structure constant \(\alpha\) is not viewed as a parameter devoid of structural origin, but rather as quantifying a granularity of surface coupling: it measures the fraction of the internal coherence budget of a proton that can effectively be expressed at its surface during its interaction with an electron. A first-order estimate, based on an active surface containing approximately \(R_{\mathrm{surf}} \simeq 500\) relations, already leads to a value close to \(\alpha^{-1}\) through the relation
\[
\alpha^{-1} \simeq \frac{R_{\mathrm{surf}}^{2}}{\mu},
\]
where \(\mu = m_p/m_e\) denotes the proton–electron mass ratio.

A refinement of this estimate reveals a more structured picture. The total number of valence stationary relations in the proton is taken to be \(r_{\mathrm{valence}} = 930\). By assigning an “average” coherence core to each valence quark, it is natural to consider one third of this quantity, \(r_{\mathrm{valence}}/3 = 310\), as a measure of the typical content associated with a local stationary regime (valence quark) within the hierarchy. Multiplication of this third by the golden ratio \(\varphi = (1 + \sqrt{5})/2 \approx 1.618\) gives
\[
R_{\mathrm{surf}}^{(\varphi)}
\simeq \varphi \,\frac{r_{\mathrm{valence}}}{3}
\approx \varphi \,\frac{930}{3}
\approx 501.59,
\]
thereby replacing the rounded estimate \(R_{\mathrm{surf}} \simeq 500\) with an expression constructed solely from quantities internal to the protonic architecture (\(930\)) and a simple geometric factor \(\varphi\).

By substituting the refined value into the relation
\[
\alpha^{-1} \simeq \frac{{R_{\mathrm{surf}}^{(\varphi)}}^{2}}{\mu}
\approx \frac{501.59^{2}}{1836.15}
\approx 137.022,
\]
with \(\mu \simeq 1836.15\), one obtains the numerical estimate \(\alpha^{-1} \simeq 137.022\). This result agrees, to within a relative deviation of a few \(10^{-4}\), with the experimental value \(\alpha^{-1} \simeq 137.036\), and is achieved without introducing any ad hoc correction terms. Consequently, the fine-structure constant can be reformulated as
\[
\alpha^{-1}
\simeq \frac{\varphi^{2}}{\mu}
\left(\frac{r_{\mathrm{valence}}}{3}\right)^{2},
\]
which expresses it exclusively in terms of the mass ratio \(\mu\), the valence relational content \(r_{\mathrm{valence}}\), and the factor \(\varphi\). 

Within this framework, the golden ratio is not invoked as an aesthetic motif or as an element of external numerology. Rather, it is proposed as a candidate for a relational geometric optimisation constant: a factor encoding the most efficient configuration for distributing the coherence content of an average stationary core toward its active surface, under the constraints of hierarchical closure and minimal leakage. In this sense, \(\varphi\) functions as a logical–topological invariant that characterises an optimal arrangement of coherence relations on the surface of stationary regimes, instead of serving as a “magic” number introduced from outside the theoretical construction. 

This interpretation is compatible with an intuition already present in the literature, according to which the fine-structure constant connects a dimensionless mass ratio \(\mu\) with a finite number of coupling channels between entities. In the present formalism, this connection is naturally recast in terms of active surfaces and relational topology. From this perspective, \(\alpha\) quantifies the fraction of action granularity, set by \(\hbar\), that remains available at the surface for coupling processes, once internal hierarchical closure has been established.

One may then reinterpret the usual relation
\[
E = \hbar \omega
\]
when applied to the intrinsic oscillation of an electron in a bound state with a proton, as the quantitative formulation of the surface-coherence expenditure required to sustain simultaneously:
\begin{itemize}[label=-]
  \item the internal stationary regime of the proton;
  \item the elementary stationary regime of the electron;
  \item and the hierarchical coupling that connects them through the active logical surface characterised by \(\alpha\).
\end{itemize}
Within this framework, \(\alpha\) is no longer a parameter without provenance. Instead, it appears as the signature of the manner in which a composite stationary regime allocates its coherence budget between internal closure and surface couplings, thereby anticipating the connection with the gravitational coupling constants introduced in the subsequent sections.

\subsection{\texorpdfstring{$k_B$}{kB} and thermal granularity}
Within this framework, the reduced Planck constant \(\hbar\) determines the minimal granularity of action associated with the pulsation cycle of an isolated stationary regime: each cycle carries one quantum of action \(\hbar\), and the proper energy of a structure is expressed as \(E = \hbar \omega\) whenever a physical realisation of this pulsation is available.

The Boltzmann constant \(k_B\) plays an analogous role as soon as one considers, not a single isolated structure, but a statistical ensemble of hierarchical structures. In conventional physics, \(k_B\) connects the mean thermal energy per degree of freedom to the temperature via relations of the form \(E_{\mathrm{th}} \sim k_B T\).[web:13] In a relational interpretation, this thermal energy may be understood as the portion of the global coherence budget that is, on average, engaged in the collective excitations of a large number of pulsation cycles.

Temperature thereby ceases to be a primitive quantity; instead, it functions as a statistical measure of the effective population of excited coherence cycles within an ensemble of structures (electrons, nucleons, atoms, etc.). The role of \(k_B\) is to relate this statistical counting of excited cycles to an average energy scale, in direct analogy with the way \(\hbar\) links an elementary cycle to a quantum of action. Thus, \(k_B\) sets the characteristic energy scale of a thermal fluctuation per degree of freedom, whereas \(\hbar\) sets the characteristic scale of an action fluctuation per cycle.[web:13]

In the present formulation of the LDP model, \(k_B\) is not yet derived in a unique and rigorous manner from the underlying axioms. A tentative expression is nevertheless proposed below, constructed solely from the internal architectures of the electron and the proton and evaluated numerically. This proposal should be regarded as a structured and testable hypothesis, rather than as a definitive derivation.

In classical thermodynamics, the Boltzmann constant \(k_B\) establishes the quantitative relationship between microscopic degrees of freedom and macroscopic thermodynamic observables, such as temperature and entropy. It appears explicitly in Boltzmann’s entropy formula \(S = k_B \ln W\) and in the energetic relation \(E \sim k_B T\), which connects an average microscopic excitation to a characteristic thermal energy scale.[web:13] From a relational perspective, \(k_B\) may be interpreted not merely as an empirical proportionality coefficient, but as quantifying a minimal granularity of coherence that emerges when a logical structure is subjected to an ensemble of possible configurations.

In this conceptual framework, the following construction examines whether the same underlying logical architecture can generate a thermal scale that is compatible, within a controlled margin of error, with the empirical value of \(k_B\).

Within the LDP (Logical/Local Dynamical Process) framework, the minimal pulsation that characterizes \(\hbar\) for the \((4,6)\) structure (associated with the electron) defines a reference period \(\tau_0\). This period is determined from the Compton wavelength of the electron, \(\lambda_C\), and the informational propagation speed \(c\), according to
\[
\tau_0 = \frac{\lambda_C}{c},
\]
where \(\lambda_C = \frac{h}{m_e c}\).[web:19] The quantity \(\tau_0\) thereby represents the logical duration of a minimal coherence cycle of an electron at rest.

For a composite structure such as the proton, the coherence energy is not confined exclusively to the \((4,6)\) stationary structure, but is instead distributed across a hierarchical architecture comprising:
\begin{itemize}[label=-]
  \item stationary sub‑structures associated with the valence quarks, characterized by \(n_u = 24\), \(n_d = 28\), \(r_u = 276\), \(r_d = 378\), yielding a total of 930 relations for the three quarks;
  \item a sea of dynamic relations \(R_{\mathrm{mer}} = 10\,087\), for a total \(R_{\mathrm{tot}} = 11\,017\), corresponding to a stationary fraction of approximately \(9\%\) and a dynamic fraction of approximately \(91\%\);
  \item an active surface comprising on the order of \(501.6\) relations, which are directly involved in the coupling with the electron and have already been employed in the construction of \(\alpha\).
\end{itemize}

Within this framework, the dynamic fraction \(R_{\mathrm{mer}}/R_{\mathrm{tot}}\) quantifies the relational degrees of freedom accessible to local thermal configurations, whereas the leakage rate \(\varepsilon\), previously introduced as the minimal coherence violation at the proton–universe interface, characterizes the fraction of coherence that is dissipated into the global relational background. 

A candidate expression for the Boltzmann constant can then be formulated by relating:
\begin{itemize}[label=-]
  \item the minimal rhythmic tension \(h/\tau_0\) associated with a single pulsation cycle of the electron,
  \item the number of active relations that participate in the logical coherence of the proton,
  \item the mass ratio \(m_p/m_e\), which determines the ratio of proper pulsation frequencies between the two structures,[web:19]
  \item the dynamic fraction of the sea and the leakage rate \(\varepsilon\), which together parameterize the fraction of coherence effectively available as thermal agitation.
\end{itemize}

Under these assumptions, one obtains an expression that involves no additional free dimensionless parameter:
\[
k_B^{\mathrm{LDP}}
= \frac{h}{\tau_0}\,
  \frac{n_{\mathrm{rel}}^{e}}{(m_p/m_e)^{2}}\,
  \frac{R_{\mathrm{mer}}}{R_{\mathrm{tot}}}\,
  \frac{\varepsilon}{4},
\]
where:
\begin{itemize}[label=-]
  \item \(h\) denotes Planck’s constant,[web:19]
  \item \(\tau_0\) is the period of the proper pulsation of an electron at rest,
  \item \(n_{\mathrm{rel}}^{e} = 6\) is the number of relations associated with the \((4,6)\) structural configuration of the electron,
  \item \(m_p/m_e = 1836.15267343\) is the proton–electron mass ratio,[web:19]
  \item \(R_{\mathrm{mer}} = 10\,087\) and \(R_{\mathrm{tot}} = 11\,017\) represent, respectively, the number of “sea” relations and the total number of relations within the proton,
  \item \(\varepsilon \approx 0.0060923\) is the leakage rate previously determined from proton coherence considerations and subsequently employed in the expressions for \(G\) and \(H_0\),
  \item the factor \(1/4\) accounts for the distribution of the leakage over the four logical groups that organize the internal structure of the nucleon.
\end{itemize}

This relation does not require the introduction of any adjustable continuous parameter. All integer quantities are fixed by the underlying relational architecture, namely \((4,6)\) for the electron and \(24, 28, 10\,087, 11\,017\) for the proton, while \(\varepsilon\) is the unique leakage rate that also appears in the formulations of the gravitational constant and the cosmological expansion rate.

By inserting the conventional numerical values \(h = 6.62607015 \times 10^{-34}\,\mathrm{J \cdot s}\), \(m_e = 9.10938356 \times 10^{-31}\,\mathrm{kg}\), \(m_p = 1.67262190 \times 10^{-27}\,\mathrm{kg}\), \(c = 2.99792458 \times 10^{8}\,\mathrm{m \cdot s^{-1}}\), and related constants into the proposed expression, one obtains
\[
k_B^{\mathrm{LDP}} \approx 1.3802 \times 10^{-23}\,\mathrm{J \cdot K^{-1}},
\]
which can be compared to the CODATA recommended value
\[
k_B^{\mathrm{exp}} = 1.380649 \times 10^{-23}\,\mathrm{J \cdot K^{-1}}.
\]
This leads to a relative discrepancy of
\[
\frac{\lvert k_B^{\mathrm{LDP}} - k_B^{\mathrm{exp}} \rvert}{k_B^{\mathrm{exp}}}
\approx 3 \times 10^{-2}.
\]

This level of agreement is modest when assessed against the precision of current experimental determinations, and the proposed formula should therefore not be interpreted as a definitive derivation of \(k_B\) within the LDP scheme. It does, however, indicate that by relying exclusively on quantities already fixed by the \((4,6)\) architecture and by the hierarchical model of proton structure (including the same leakage parameter \(\varepsilon\) that appears in the expressions for \(G\) and \(H_0\)), one can relate classical thermodynamic discreteness to the same relational framework that underlies mass, electric charge, and gravitation. Such a convergence, even if only approximate, supports the hypothesis that macroscopic thermodynamics may emerge as a statistical manifestation of logical coherence constraints operating already at the level of nucleonic organization. 

The proposed expression should thus be regarded as conjectural. It depends on the identification of the integers \(6\), \(10\,087\), \(276\), and \(378\) with specific numbers of relations in the protonic architecture, and on the reuse of the same leakage parameter \(\varepsilon\) that is employed for \(G\) and \(H_0\). Nonetheless, the fact that a parameter–free combination determined solely by the underlying logical structure (i.e., without introducing any additional continuous free parameter) yields a value of \(k_B\) with a relative error of order \(3 \times 10^{-2}\) lends support to the broader proposal that macroscopic thermodynamics is linked to the same relational substratum as microphysics and cosmology.

\section{Protons, Black Holes, and Surface-Encoded Information}

In contemporary physics, two ostensibly disparate entities occupy opposite extremes of the characteristic mass and length scales: the proton, a fundamental constituent of ordinary baryonic matter, and the black hole, the archetypal manifestation of strong-field gravitation. Nevertheless, both systems can be understood as maximally saturated information-bearing configurations, in which the physically relevant degrees of freedom are effectively encoded on a finite boundary surface rather than distributed throughout an undifferentiated bulk volume.

\subsection{The proton as a hierarchical information core}
Within the modern theoretical description, the proton cannot be regarded as a simple triplet of pointlike quarks. Quantum chromodynamics (QCD) establishes that its mass and energy are generated predominantly by the non‑perturbative dynamics of the gluon field and the “sea” of virtual quark–antiquark pairs, rather than by the bare masses of the three valence quarks. The proton’s internal structure therefore exhibits features akin to a fractal‑like branching of dynamical degrees of freedom, in which the valence quarks constitute only the most accessible manifestation of a substantially denser underlying organisation. 

In the relational framework adopted here, this internal complexity is reinterpreted as a hierarchical architecture whose finite, dynamically active boundary concentrates the effective coupling channels with the electron. Within this view, certain dimensionless fundamental constants can be reformulated as quantitative measures of coherence between the proton’s internal hierarchical structure and its interaction surface.

\subsection{Black holes: the event horizon as an entropy carrier}
On the gravitational side, the pioneering work of Bekenstein and Hawking established that black holes cannot be regarded as mere spacetime regions from which no signal can escape. In order for the second law of thermodynamics to remain valid when matter crosses the event horizon, it is necessary to attribute to black holes an entropy
\[
S_{\mathrm{BH}} = \frac{k c^{3} A}{4 G \hbar},
\]
which is proportional to the area \(A\) of the event horizon. This entropy is naturally interpreted as an information-theoretic entropy: in accordance with the Bekenstein bound and the holographic principle, the horizon encodes the maximal number of bits that can be contained within a given spatial region.

Hawking subsequently demonstrated that black holes emit thermal radiation, thereby determining the temperature associated with the event horizon and confirming that these thermodynamic quantities are not merely formal constructs. Rather, black holes obey bona fide laws of thermodynamics, with the horizon area playing the role of entropy. In several contemporary frameworks—such as the membrane paradigm, “horizon-as-apparatus” models, and holographic dualities—the relevant microscopic degrees of freedom are attributed primarily to the horizon surface rather than to the black hole interior. The largest known gravitationally bound objects thus emerge as systems whose physically relevant informational content is effectively localized on a finite-area surface, measured in units of the Planck area.

\subsection{An inverted mini black hole}

On the basis of the two preceding descriptions, the analogy can be formulated more systematically. In both cases, one observes convergence towards the same organisational pattern: an extremely dense configuration of degrees of freedom whose internal complexity is, for an external observer, effectively summarised by a finite set of surface parameters. The correspondence can be articulated as follows:
\begin{itemize}[label=-]
    \item In the case of an astrophysical black hole, the event horizon encodes the total accessible entropy: the information concerning all matter and radiation that have crossed the horizon is, in an effective sense, “stored” on this surface, and the exterior geometry is determined solely by a small number of global charges (mass, electric charge, angular momentum).
    \item In the case of the relational proton, the internal hierarchical architecture (saturated valence-quark sector and dynamically structured relational sea) fixes a global coherence budget; however, only a finite “active surface” effectively participates in external couplings. It is this relational interface, or “skin”, that is probed by dimensionless observables such as the fine-structure constant \(\alpha\) or mass ratios.
\end{itemize}

In this perspective, one may characterise the proton as an inverted mini black hole. Whereas an astrophysical black hole concentrates the macroscopic information of a spacetime region on a geometric horizon, the proton concentrates the information of a minimal composite structure on an internal relational closure, whose active logical surface plays, with respect to the electron and other quantum fields, a role that is functionally analogous to that of a horizon. In both situations, the observable external dynamics are constrained primarily by the information structure accessible on a finite surface, rather than by the detailed properties of the internal “volume”.

\subsection{It from bit: Hawking, Bekenstein, and Wheeler}
Wheeler encapsulated this convergence in the now canonical expression “it from bit”: the hypothesis that each constituent of the physical world ultimately emerges from elementary binary information units—yes/no alternatives—rather than from a pre-given continuous substrate. In the specific case of black holes, this research programme attains a quantitative formulation through the Bekenstein–Hawking entropy: the event horizon is endowed with an informational entropy that must be regarded as a bona fide physical quantity.

The relational framework proposed here transfers this intuition to the hadronic sector. The proton is represented as a combinatorial organisation of agreement/disagreement relations, a subset of which is frozen into locally stationary configurations (valence quarks), while another subset sustains a collective dynamical regime (the relational sea). The active surface of the proton, defined as the minimal relational envelope through which this structure can couple to an electron described by \((4,6)\), plays a dual role: on the one hand, it compresses internal coherence into a finite set of coupling channels; on the other, it permits the interpretation of dimensionless constants, such as the fine-structure constant, as ratios of internal/surface coherence rather than as primitive, unexplained parameters.

From this perspective, the “it from bit” thesis is instantiated at two distinct scales: at the macroscopic level, through Bekenstein–Hawking black holes, and at the microscopic level, through the proton conceived as a minimal composite relational structure. In both regimes, a finite organisation of information, localised on a surface, acquires an ontological status equivalent to that of a physical object: the event horizon in the case of the black hole, the active surface in the case of the proton. The analogy does not equate the proton with a black hole in the relativistic sense; rather, it emphasises that, from the standpoint of coherence and information, the smallest composite constituent of matter and the largest known gravitational structure obey an analogous principle: physical reality becomes manifest where an informational structure attains a saturation threshold that renders it irreducible to a purely linguistic or descriptive construct.

\section{From coexistence to an emergent metric}
The relational characterisation of closed structures, together with the interpretation of proper time as the discrete counting of pulsation cycles, provides an adequate description of an isolated hierarchical structure. Empirical physics, however, is not confined to single, isolated entities: electrons, protons, and neutrons coexist in large ensembles and combine into nuclei, atoms, and macroscopic bodies, all while preserving their individual identities.

Within the logical framework, the coexistence of multiple hierarchical structures entails that their respective stationary regimes must be capable of maintaining internal coherence without mutually destabilising or annihilating one another. This condition does not yet constitute an interaction in the dynamical sense; rather, it specifies a compatibility constraint: certain configurations of hierarchical closures can occupy a common “logical neighbourhood”, whereas others must be mutually excluded in order to preclude logical inconsistency or contradiction.

From these compatibility relations there emerges a network of neighbourhoods: each hierarchical structure is directly connected to a finite set of other structures with which it is compatible, as determined by charge conditions and coherence constraints. This network presupposes neither continuity nor dimensionality nor any pre-given geometric structure; it encodes solely the conditions for the possibility of stable coexistence.

In the case of composite structures such as protons, this compatibility is directly influenced by the minimal coherence leakage associated with each hierarchical level. The fraction \(\varepsilon\) of coherence that does not fully close within a single proton must necessarily be absorbed by the surrounding network: it appears as an additional constraint on the manner in which hierarchical closures can occupy a common logical neighbourhood. Any configuration that locally increases the density of protons thereby alters the leakage pattern and enforces a readjustment of compatibilities in its immediate vicinity. 

Within this network, the notion of neighbourhood defines a rudimentary topology: two configurations are said to be neighbours if their respective closures can be sustained simultaneously without violating the coherence constraints. The resulting “space” is thus not a passive background, but a collective property emerging from the organisation of coexistence relations among stationary structures. 

On this basis, the minimal metric introduced earlier acquires a precise operational meaning: it supplies a criterion for comparing separations between hierarchical structures, not in terms of pre-given distances, but in terms of the differential coherence cost required to maintain them within the same logical neighbourhood. One configuration will be considered “farther” than another if its coexistence with all already present structures demands a larger readjustment of the coherence distribution. 

This transition from a purely binary notion of compatibility to a graded comparison of separations constitutes a first step towards an emergent metric in the physical sense. It underpins an interpretation of gravitation not as an external force superimposed upon a pre-established space, but as the systematic manifestation of these coherence readjustments when the distribution of hierarchical structures is modified. 

It is crucial to stress that this emergent metric does not arise at the level of an isolated minimal structure of type \((4,6)\). As long as only a single elementary stationary pulsation is realised, there exists neither a neighbourhood network nor a requirement to compare separations: geometry, in the physical sense, is not yet defined. Only once a threshold is crossed—namely, when composite hierarchical structures coexist in sufficient number, each characterised by a minimal leakage \(\varepsilon\)—does the cost of coherence readjustments become a collective quantity, interpretable as an effective curvature of the emergent metric. Gravitation can then be understood as the quantitative expression of these readjustments, as will be made explicit in the following section.

\section{Gravitation as Variation of Coherence Cost}

A closed hierarchical structure, such as an electron or a nucleon, is characterised by an internal coherence cost, quantified both by the number of relational constraints necessary for its maintenance and by the intrinsic pulsation associated with its rest mass. This cost is not superimposed on a neutral background substrate; rather, it represents the fraction of the global coherence budget that the relational framework must allocate to this structure in order to preserve its stability.

When multiple hierarchical structures coexist, their coherence costs do not simply sum within an otherwise indifferent background. Each structure imposes specific constraints on the admissible modes of coherence distribution in its logical neighbourhood. As the density of closed hierarchical structures in a given region increases, the allocation of coherence required to sustain all stationary regimes simultaneously becomes more constrained.

Consider, for example, two nucleons initially separated by a large distance with respect to the emergent metric. To first order, their internal closures may be treated as approximately independent, and the coherence readjustments demanded of the relational framework remain minimal. As these nucleons are brought into closer proximity in the network of logical neighbourhoods, the mutual compatibility of their stationary regimes necessitates a more substantial reorganisation of the coherence budget within the region they jointly occupy. The total coherence cost attributed to this common region accordingly increases.

This increase does not correspond to an external injection of energy; rather, it reflects a redistribution of the coherence already present. The relational framework must concentrate more of its resources to maintain, within the same spatial domain, several hierarchical closures that remain mutually compatible. In a metric description, this concentration of coherence manifests itself as an effective curvature: the trajectories that minimise coherence readjustments cease to be geodesics of a flat background and instead bend around the regions where the coherence cost attains its highest values.

The minimal coherence leakage associated with each proton then becomes a determining factor. As long as closure is exact, as in the case of an isolated \((4,6)\) elementary structure, no constraint is imposed on the emergent metric: all coherence circulation remains confined within the structure itself. Once a composite hierarchy exhibits a nonzero leakage \(\varepsilon\), a fraction of this circulation must be supported by the surrounding relational framework. Gravitation can then be interpreted as the collective manifestation of these minimal leakages, when a large number of nucleons simultaneously impose their constraints on the network of logical neighbourhoods. 

In this perspective, gravitation can be understood as the tendency of hierarchical structures to follow trajectories of minimal coherence cost. These structures do not “attract” one another in the sense of a Newtonian force; rather, they evolve within a framework in which the emergent metric is continuously adjusted according to the distribution of coherence costs, and their trajectories align with the directions that minimise these adjustments. 

Within this framework, what is conventionally termed “gravitational mass” corresponds to a quantitative measure of the hierarchical coherence cost of a closed structure, while the curvature of the emergent metric encodes the spatial distribution of this cost within the network of logical neighbourhoods. Gravitation no longer appears as a field superimposed on a pre-existing space; instead, it emerges as the large-scale organisation of the collective redistribution of coherence necessitated by the coexistence of persistent hierarchical structures.

\section{\texorpdfstring{Leakage parameter \(\varepsilon\) and hierarchical amplification}{Leakage parameter epsilon and hierarchical amplification}}

The representation of the proton as a minimal composite structure reveals the existence of a coherence leakage: a fraction of the relational circulation required for its stability does not close exclusively within its internal network of relations, but partially transits through the surrounding relational framework. This leakage does not correspond to an energetic loss; rather, it reflects the structural constraint that a hierarchy cannot simultaneously realise maximal closure at all subordinate levels and an absolutely isolated closure at the global level.

To quantify this controlled deviation from perfect closure, one introduces a leakage parameter \(\varepsilon\). At the proton scale, \(\varepsilon\) characterises the relative discrepancy between an ideal hierarchical closure—where all internal relations would close exactly within the composite structure—and the effective closure, in which a small fraction of coherence is allowed to project outward. The parameter \(\varepsilon\) is therefore a dimensionless quantity intrinsically associated with the composite hierarchy itself, and not defined with respect to any pre‑imposed metric background.

This local leakage is not restricted to the hadronic scale. When a large number of protons (or nuclei) assemble into higher‑order structures (atoms, solids, planets, stars), the coherence leakages associated with each constituent do not merely sum linearly; instead, they combine and undergo hierarchical amplification. At each organisational level, a small fraction of internal coherence escapes perfect closure and is transmitted to the next level, such that the cumulative effect can become significant at macroscopic scales and manifest as a deformation of the emergent metric.

Within this framework, \(\varepsilon\) acquires a central role in a reinterpretation of gravitation. At the scale of an isolated proton, its magnitude is extremely small and produces no directly observable deviation from standard behaviour. However, when one examines the propagation of this leakage across multiple hierarchical levels (from the internal structure of the proton through ordinary matter up to astrophysical configurations), the same parameter appears within an effective exponent that governs the dilution, or rarefaction, of coherence at large scales. It is this exponential dependence that will enter into the proposed expression for an effective gravitational coupling constant.

At the current stage of development of the model, it is essential to emphasise that the numerical value of \(\varepsilon\) has not yet been obtained by a purely combinatorial derivation. Instead, it is fixed by imposing that the resulting expression for the gravitational coupling constant reproduces the empirically measured value of Newton’s gravitational constant. The parameter \(\varepsilon\) must therefore be regarded as observationally constrained rather than as a quantity fully deduced from the underlying axioms. Any subsequent derivation of \(\varepsilon\) within the theory will have to recover this empirically determined value as a consequence of the proposed hierarchical structure; otherwise, the reinterpretation advanced here would remain theoretically incomplete.

In the following section, this leakage parameter is employed, together with the internal structure of the proton and the fine-structure constant, to construct an explicit expression for the gravitational coupling constant. The central issue is then to assess to what extent the relational framework can account for the observed order of magnitude of gravitation without the introduction of an additional, independent fundamental constant.

\section{Towards a gravitational coupling constant}
If gravitation is interpreted as the manifestation of spatial and temporal variations in the collective coherence cost, one may introduce a gravitational coupling constant as a quantitative measure of the sensitivity of the emergent metric to this cost. Such a constant then establishes a quantitative relationship between the hierarchical content of a closed structure and the curvature it induces in its surrounding region. 

Each nucleon can be characterised by three fundamental quantities:
\begin{itemize}[label=-]
  \item a total number of internal relations, which quantifies its hierarchical coherence cost;
  \item a Compton wavelength associated with its rest mass, which defines a minimal spatial closure scale;
  \item a Compton frequency, which determines the rate at which its internal coherence is re‑established and thus specifies its elementary proper time.
\end{itemize}
These quantities permit the definition of a coherence density associated with a nucleon, understood as a relational cost per minimal closure volume and per elementary proper oscillation cycle. In a configuration in which the nucleon density is spatially non‑uniform, the corresponding coherence density is likewise position‑dependent, thereby enforcing a reorganisation of the emergent metric. An increase in coherence density within a given region is then associated with an increase in metric curvature, in the sense that the geodesics—i.e., the paths of minimal coherence cost—are correspondingly deflected.

The gravitational coupling constant can now be specified in a more precise and structurally motivated manner. It must connect the variation of coherence density—amplified by the minimal leakage \(\varepsilon\) associated with each proton—to the effective curvature of the emergent metric. A natural parametrization is obtained by combining:
\begin{itemize}[label=-]
  \item the elementary action \(\hbar\), already interpreted as the minimal quantum of action associated with an elementary logical stationary regime;
  \item the limiting speed \(c\), which sets the fundamental scale for information transfer. Within the relational framework, this limiting speed governs not only electromagnetic propagation, but also the maximal rate at which coherence readjustments can propagate in the emergent metric, thereby defining the upper bound of any causal influence;
  \item the proton mass \(m_p\), which encodes its global hierarchical coherence cost;
  \item the fine‑structure constant \(\alpha\), which characterizes the strength of active surface interactions. In the LDP model, these active surface interactions correspond to a finite set of logical links that remain available for external couplings once the internal coherence of the proton has been saturated. A detailed combinatorial analysis—based on the valence‑quark / quark‑sea hierarchy and on the geometry of closure—leads to an estimate of an active surface comprising approximately 502 relations, which allows \(\alpha\) to be reinterpreted as a surface‑coherence ratio;
  \item the leakage parameter \(\varepsilon\), which quantifies the minimal logical violation of hierarchical closure.
\end{itemize}

Once the functional roles of \(\hbar\) and \(c\) have been identified at the level of elementary stationary regimes, and the structuring role of \(\alpha\) in surface interactions has been clarified, it becomes possible to examine how a gravitational coupling constant may emerge from the imperfect hierarchical closure of hadronic systems.

A possible expression for the gravitational coupling constant is therefore
\[
G_{\mathrm{LDP}}
= \frac{2 \hbar c}{m_p^{2}}\,
  \frac{1}{\alpha^{18(1-\epsilon)}},
\]
where the exponent \(18(1-\epsilon)\) encodes the minimal number of fractal (hierarchical) iterations required for the deformation induced by the leakage to propagate from the proton scale to macroscopic scales. This iteration count is inferred from a detailed hierarchical model of proton structure. Its exact combinatorial derivation is not provided in the present work and will be developed explicitly in a separate, dedicated analysis.

This relation does not arise from an unconstrained fit to the experimental value of \(G\); rather, it emerges from a restricted combination of the physical elements emphasized in the preceding sections: \(\hbar\) represents the quantum of action per cycle, \(c\) the maximal signal velocity governing the propagation of coherence readjustments, \(m_p\) the hierarchical coherence cost associated with the proton, \(\alpha\) the fraction of surface coherence effectively available for couplings, and \(\epsilon\) the minimal leakage induced by the composite hierarchical structure. The numerical factor 18 characterizes the minimal number of hierarchical steps required for proton-level leakage effects to percolate to macroscopic scales. Its detailed combinatorial origin is beyond the scope of this article and will be specified in a subsequent technical work.

The factor 18 is thus not introduced as an adjustable parameter. It corresponds to the minimal number of hierarchical iterations needed for an elementary coherence perturbation—initiated by the leakage \(\epsilon\) at the proton scale—to propagate iteratively up to macroscopic domains. This number results from counting the groups of active internal relations that contribute to the hierarchical closure of the proton (including valence quarks, the quark sea, and the active surface), as formulated in the complete version of the model.

The detailed counting procedure, and the manner in which it leads to 18 effective iterations, are presented in the full exposition of the LDP model, to which the present text refers for all combinatorial aspects.

By inserting the measured values of \(\hbar\), \(c\), \(m_p\), and \(\alpha\), and adopting \(\varepsilon \approx 0.0060923\), one obtains the numerical estimate:
\[
G_{\mathrm{LDP}} \simeq 6.6743 \times 10^{-11}\,\mathrm{m^{3}\,kg^{-1}\,s^{-2}}.
\]

The construction of \(G_{\mathrm{LDP}}\) must be interpreted with particular caution. It does not yet follow from a full derivation grounded solely in the axioms of the model; rather, it is obtained from a constrained combination of previously introduced quantities: \(\hbar\) as the quantum of action per cycle, \(c\) as the upper bound on the propagation speed of coherence readjustments, \(m_p\) as the hierarchical coherence cost associated with the proton, \(\alpha\) as a measure of surface coherence, and \(\varepsilon\) as the minimal leakage parameter characterising the composite hierarchy.

The exponential factor \(18(1-\varepsilon)\) encodes the hypothesis that a microscopic leakage, introduced at the proton scale, propagates through a finite number of hierarchical iterations before becoming observable at macroscopic scales. In the present state of development of the model, however, the number 18 is not yet supported by an explicit combinatorial derivation: it effectively summarises a provisional and still incomplete hierarchical scheme, which will need to be made explicit in a separate, more technical treatment.

Likewise, the numerical value of \(\varepsilon\) employed in the expression of \(G_{\mathrm{LDP}}\) is not presently predicted in an autonomous manner. It is chosen such that \(G_{\mathrm{LDP}}\) reproduces the empirically measured gravitational constant, which makes \(\varepsilon\) an observationally calibrated parameter rather than one derived from the internal structure of the proton. Any future derivation will have to recover this value as a consequence of the axioms and of the hierarchical construction; failing this, the current expression can only be regarded as a consistent reinterpretation of the data, rather than a prediction obtained strictly from first principles.

In summary, \(G_{\mathrm{LDP}}\) provides an interesting consistency test: it shows that a gravitational coupling constant can be reconstructed from the relational elements already introduced, with the correct order of magnitude, but it does not yet constitute a complete derivation nor an independent prediction of gravitation. 

At the present stage, the numerical value of \(\epsilon\) has not yet been obtained through a purely combinatorial derivation. Operationally, it is fixed by requiring that the expression for \(G_{\mathrm{LDP}}\) reproduce the experimentally measured value of Newton’s gravitational constant. This calibration should be interpreted as a consistency check of the model: it constrains the admissible range of \(\epsilon\) and sets a quantitative target that any subsequent microscopic derivation will be required to recover. 

The relative discrepancy is of order \(10^{-4}\), which is compatible with current experimental uncertainties on \(G\) and with the approximations inherent in the hierarchical description of the proton. This agreement does not constitute a definitive validation of the model, but it does demonstrate that a gravitational coupling constant can emerge quantitatively from:
\begin{itemize}[label=-]
  \item the internal structure of the proton;
  \item the distribution of its active relational degrees of freedom;
  \item and a minimal coherence leakage \(\epsilon\);
\end{itemize}
without postulating a fundamental gravitational interaction independent of the relational hierarchy. 

In principle, a gravitational coupling constant may be defined as the proportionality factor between:
\begin{itemize}[label=-]
  \item the variation in coherence density induced by the introduction of an elementary hierarchical structure (for instance, a proton) within a given region;
  \item and the variation in the effective curvature of the emergent metric, as operationally determined by the deviation of the minimal‑coherence trajectories in that region.
\end{itemize}
In a formulation closer to established theories, this constant plays an analogous role to that of Newton’s gravitational constant: it relates a quantity associated with material content (here, the hierarchical coherence cost) to a measurable deformation of the metric. The crucial difference is that, within the relational framework, this constant is not postulated as a fundamental parameter; rather, it emerges from the internal structure of nucleons and from the manner in which their hierarchical closures combine within the network of logical neighbourhoods.

A coherent numerical estimate can, in principle, be obtained by combining the total number of internal relations characterizing a nucleon, the radius associated with its Compton wavelength, and its proper frequency, in such a way as to construct a quantity with the physical dimensions of a gravitational coupling constant. The fact that such a construction yields a value of the correct order of magnitude, compatible with the experimentally measured gravitational constant, suggests that interpreting the hierarchical coherence cost in gravitational terms is not only conceptually admissible but also quantitatively credible.

\subsection{Status of the model’s internal parameters}
The foregoing construction involves several internal numerical parameters. Before concluding, it is necessary to clarify their epistemological status, since they do not play equivalent roles within the model:
\begin{itemize}[label=-]
  \item the leakage parameter \(\epsilon\);
  \item the hierarchical factor 18 in the exponent of \(\alpha\);
  \item the integers associated with local architectures \((4, 6, 24, 28, 500, 10\,087, 11\,017, \text{etc.})\).
\end{itemize}

Some of these numerical values follow directly from the axioms and from standard results in graph theory:
\begin{itemize}[label=-]
  \item the minimal \((4,6)\) structure is implied by the constraints of reproducibility, stability of reference, logical autonomy, and absence of an a priori hierarchy;
  \item the expression \(r(n) = n(n-1)/2\) is a consequence of modelling local stationary regimes as complete graphs;
  \item the requirement of a non-zero coherence leakage in composite structures (nucleons) follows from the impossibility of simultaneously achieving maximal local and global closure.
\end{itemize}

Other integer values are introduced as minimal solutions to specific combinatorial constraints, without having yet been proven unique:
\begin{itemize}[label=-]
  \item the values \(n_u = 24\) and \(n_d = 28\) for the valence \(u\) and \(d\) quarks;
  \item the effective number of surface relations \(R_{\mathrm{surf}} \simeq 502\);
  \item the total number of internal relations of the proton \(r_{\mathrm{total}} = 11\,017\).
\end{itemize}
These choices are not arbitrary: they arise from a systematic exploration of architectures consistent with the axioms (minimal completeness, minimal leakage, stationary fraction \(\sim 9\%\), dynamic fraction \(\sim 91\%\)), and they successfully reproduce several empirical ratios (for instance, the predominance of dynamic energy in the proton mass and the value of \(\alpha\) with an accuracy better than \(1\%\)). They should nevertheless be regarded as strongly constrained conjectures, which may be refined or more rigorously justified in a more detailed combinatorial analysis.

Finally, several parameters are not yet obtained through purely deductive reasoning, but are instead determined by calibration against experimentally measured constants:
\begin{itemize}[label=-]
  \item the numerical value \(\epsilon \simeq 0.0060923\) is chosen such that the effective gravitational coupling \(G_{\mathrm{LDP}}\) matches the empirically measured Newtonian gravitational constant;
  \item the factor 18 appearing in the exponent of \(\alpha\) represents the minimal number of hierarchical iterations required for the propagation of coherence leakage from the proton scale to macroscopic scales; its detailed combinatorial origin is left for a subsequent, more explicit development.
\end{itemize}

In the current formulation of the model, these parameters are therefore only semi‑fundamental:
\begin{itemize}[label=-]
  \item their qualitative necessity (existence of a leakage parameter \(\epsilon > 0\), existence of a hierarchical amplification factor) is enforced by the internal logical structure of the framework;
  \item their precise numerical values are constrained by the requirement of compatibility with observed physical constants (in particular with \(G\)).
\end{itemize}

A fully empirical and theoretical completion of the model will need to demonstrate how these values emerge from a detailed combinatorial enumeration of active relations within hadronic structures, without appealing to empirical calibration. In the present work, the calibration procedure plays the role of a consistency check: it establishes that a single relational framework can reproduce, on the basis of an internal nucleon architecture, both the fine-structure constant \(\alpha\) and an effective gravitational coupling \(G_{\mathrm{LDP}}\) that is numerically consistent with the measured value of \(G\). In summary, the LDP model is organized on three conceptual levels: a stringent deductive logic as its foundation, a conjectural yet tightly constrained relational architecture at its core, and an empirical calibration that furnishes a precise quantitative target for any prospective complete theory.

\subsection{Sketch of the 18 coherence filters and residual gravitation}

The proposed expression for the effective gravitational coupling constant,
\[
G_{\mathrm{LDP}} \simeq \frac{2 \,\hbar c}{m_{p}^{2}} \,\frac{1}{\alpha^{18(1-\varepsilon)}}
\]
contains a discrete exponent \(18\), multiplied by the leakage factor \((1-\varepsilon)\). In the foregoing analysis, this integer was interpreted as the “minimal number of hierarchical iterations” necessary for the coherence leakage associated with a proton to emerge, on macroscopic scales, as an observable deformation of the emergent metric. 

In this section, a more detailed schematic interpretation of these eighteen filtering stages is proposed. Specifically, they are related to the constraints already introduced in the construction of the \((4,6)\) building block, in the definition of the \(u/d\) valence blocks, and in the hierarchical internal architecture of the proton. Within this framework, each “filter” is to be understood as a distinct level of structural, dynamical, or symmetry-based constraint that selectively attenuates or redistributes coherence, such that only a small residual component survives to contribute effectively to long-range gravitational interaction.

Within the LDP framework, coherence is not conceived as being continuously distributed over a neutral background, but rather as circulating in the form of discrete “packets” within relational structures that exhibit a tendency toward closure. Each time a coherence packet traverses a closure filter, a portion of its capacity for outward projection is expended in favour of increased internal stability or of shorter-range couplings (for example, electromagnetic interactions). Gravitation may then be interpreted as a collateral effect of this multistage process: after a finite number of such filterings, what remains is an extremely diffuse residual constraint, insufficient to locally organize matter but, when considered in the aggregate, adequate to induce curvature in the emergent metric.

From this standpoint, the exponent \(18\) is understood to represent the number of significant coherence filters that the packets must cross between the scale of the minimal \((4,6)\) brick and the scale of macroscopic aggregates at which the gravitational constant \(G\) is operationally determined. Without asserting the uniqueness of this factorization, one may nonetheless propose the following grouping, based solely on the constraints already specified.

These conditions specify the very possibility of a minimal stationary regime (identified with the electron):
\begin{itemize}[label=-]
    \item Existence of an intrinsic binary pulsation (a logical 2-cycle) acting on the edge signs.
    \item Local triangular coherence, i.e., the product of edge signs is \(+1\) for every existing triangle.
    \item Completeness of the \((4,6)\) graph, in the sense that every pair of vertices is mutually connected.
    \item Connectivity of the block of coherent triangles (presence of at least two triangles sharing an edge, together covering the entire structure).
    \item Structural equivalence of the vertices (absence of any internal hierarchy within the \((4,6)\) configuration).
    \item Minimality: no proper subgraph simultaneously satisfies all of the above properties.
\end{itemize}
Taken together, these six criteria ensure that coherence remains strictly confined within a single brick, with no leakage to the exterior.

The organisation of multiple bricks into a composite hierarchical structure (identified with the proton), while preserving global stability, is governed by the following conditions:
\begin{itemize}[label=-]
    \item The number of entities associated with each quark is a multiple of 4, corresponding to a superposition of logical tetrads.
    \item Each quark satisfies an internal closure constraint, represented by a complete graph \(K_{n_{q}}\) (quasi-saturated internal triangular coherence).
    \item The cardinalities \(n_{u} = 24\) and \(n_{d} = 28\) are chosen as the first pair of values that simultaneously fulfil completeness, minimise asymmetry, and maximise the optimisation functional \(\Omega\).
    \item The global electric charge \(+1\) is implemented via the internal configuration \((u,u,d)\), in a manner compatible with the charge \(-1\) attributed to the \((4,6)\) brick.
    \item The relational content is decomposed into a stationary core (930 relations) and a dynamical sea (10\,087 relations), with a relative proportion of approximately \(9\%\) to \(91\%\), respectively.
    \item The sea displays sufficient connectivity to mediate coherence transfer between valence blocks, while remaining below the threshold of global saturation that would otherwise disrupt the hierarchical organisation.
    \item A finite active surface exists, strictly contained within the set of valence relations, which is capable of establishing logical coupling to the electron via coherent mixed triangles.
    \item The global closure of the proton is preserved despite this internal subdivision into valence blocks, sea, and active surface.
\end{itemize}
These eight constraints characterise the redistribution of coherence packets within the proton: one fraction is confined within the local closure structures associated with the quarks, another fraction is transferred to the dynamical sea, and a further fraction remains accessible at the surface for external couplings.

In addition, four supplementary filtering mechanisms regulate the transition through which the residual coherence packets depart from the strictly internal domain and become manifest within the surrounding relational framework:
\begin{itemize}[label=-]
    \item Logical proton–electron compatibility: the existence of coherent mixed triangles that implement a \(+1/-1\) charge coupling without triggering an uncontrolled proliferation of frustrated or violated triangles.
    \item Attenuation of the fraction of coherence accessible via surface couplings, parametrised by the fine-structure constant \(\alpha\), which sets the maximal “resolution” or granularity at which such exchanges can occur.
    \item Logical impossibility of achieving perfect closure simultaneously at all structural levels (bricks, \(u/d\) blocks, sea, surface), which enforces the existence of a minimal leakage \(\varepsilon > 0\) toward configurations external to the 11\,017 internal relations.
    \item Hierarchical propagation of this leakage across successively higher aggregation scales (nucleons, nuclei, atoms, ordinary matter, astrophysical structures), such that, within a metric description, it ultimately appears as an effective curvature encoded by \(G_{\mathrm{LDP}}\).
\end{itemize}

Taken together, these four filters provide an explicit description of the transfer of coherence packets from the strictly internal domain (proton + electrons) to the collective relational framework, and subsequently to macroscopic scales. At each stage, the direct contribution of coherence is either absorbed by electromagnetic couplings (through \(\alpha\) and the active surface) or by the internal stabilisation of the hierarchical structure; only an extremely small residual component survives as a background constraint, which is identified with the gravitational contribution.

The foregoing characterisation of the eighteen filters must be understood as a structured conjecture rather than as the outcome of a fully completed combinatorial derivation: in a more refined analysis, several filters may prove to be logically dependent, or may have to be decomposed into more elementary sub-constraints. Nonetheless, this framework indicates that the exponent \(18\) can be interpreted not as an arbitrary integer but as a discrete summary of the number of layers of coherence constraints that coherence packets must traverse for gravitation to emerge as a residual effect of the hierarchical closure of forms.

Future work will need to specify, in a purely combinatorial formalism (signed graphs, classes of triangles, orbits of the internal dynamics), how these filters correspond to rigorously defined invariants, and to demonstrate in what sense the value \(18\) is minimal for simultaneously reproducing the internal structure of the proton, the empirical value of \(\alpha\), and the observed order of magnitude of \(G\).

\section{Cosmological Constant and Residual Coherence}

Within the framework of general relativity, the cosmological constant \(\Lambda\) enters Einstein’s field equations as an additional term that can be interpreted as a homogeneous vacuum energy density responsible for the accelerated expansion of the Universe. Contemporary cosmological observations indicate that this component, commonly identified with dark energy, provides the dominant contribution to the total energy budget of the Universe, while simultaneously exhibiting an extremely small numerical value in natural units. The tension between this observed smallness and naïve estimates of quantum vacuum energy constitutes the well-known “cosmological constant problem”.

In the LDP framework, space–time is not treated as a fundamental entity; rather, it is regarded as emergent from the compatibility of hierarchical closures and from the network of logical neighbourhoods they engender. Gravitation is reinterpreted as the trace, in an effective metric description, of a minimal coherence leakage \(\varepsilon\), which is amplified across several hierarchical levels. From this standpoint, it is natural to consider whether the cosmological constant might be associated with a minimal density of residual coherence that remains distributed throughout the global relational network, outside any explicitly identified hierarchical structures (such as nucleons, atoms, or galaxies).

One may thus conjecture that \(\Lambda\) encodes, at the level of the effective large-scale geometry, the presence of a coherence background that is neither captured by local closures (particles, nuclei) nor by intermediate closures (stars, galaxies), but instead persists as a residual “tension” of the network on cosmological scales. This coherence density of the “logical vacuum” would then manifest itself, in the metric description, as a non-zero mean curvature even in regions devoid of ordinary matter, thereby reproducing the phenomenological role of a positive cosmological constant in \(\Lambda\)CDM-type models.

At the current stage of development, this proposal remains qualitative. No quantitative expression for \(\Lambda\) has yet been derived: the LDP framework does not determine the numerical value of the cosmological constant and does not resolve the established hierarchy problem between the vacuum energy predicted by quantum field theories and the observed value of \(\Lambda\). The claim is restricted to suggesting that, in a setting where space–time is emergent, the effective cosmological constant may be reinterpreted as a measure of the large-scale residual coherence of the relational network, rather than as an additional fundamental parameter.

The same leakage parameter \(\varepsilon\), originally introduced to achieve closure of the proton’s internal structure, also appears in the expressions for the gravitational constant \(G\) and subsequently for the Hubble parameter \(H_0\), yielding \(H_0 \approx 64.7\,\mathrm{km\,s^{-1}\,Mpc^{-1}}\), a value compatible with current cosmological constraints. The use of a single leakage rate, consistently applied from the nucleonic scale up to cosmological distances, is one of the most distinctive features of the framework, as it links a local coherence condition to a globally extensive cosmological quantity.

\section{Bound systems and effective gravitation}
Once introduced, the gravitational coupling constant enables, at an effective level, the description of how ensembles of nucleons behave as collective sources of curvature for the emergent metric. Composite systems, ranging from atomic nuclei to macroscopic bodies, can then be interpreted as hierarchical assemblages of elementary coherence costs. 

The gravitational mass of a finite system is not given by the mere sum of its rest masses considered in isolation. Each level of aggregation modifies the distribution of coherence costs: certain contributions add constructively, whereas others partially cancel, depending on the organisation of the hierarchical closures. The gravitational coupling constant parametrises the discrepancy between this naïve additive estimate of costs and the collective effect that is observable in the emergent metric. 

Within the relational framework, the gravitational “force” experienced by a given structure constitutes only a coarse‑grained description of the deviation of its minimal‑coherence trajectory when it propagates through a region with inhomogeneous coherence density. The law of free fall can thus be reformulated as follows: for a fixed coupling, all structures locally follow the same minimal‑cost trajectories, independently of their internal coherence cost, provided that the latter remains negligible in comparison with that of the surrounding medium.

\section{Ordinary matter and hierarchical organisation}
Ordinary matter may be interpreted as the outcome of a repeated hierarchical organisation of closed structures: electrons, nucleons, atoms, molecules and macroscopic aggregates. Each of these organisational levels corresponds to a distinct mode of allocating a finite “coherence budget” between quasi‑stationary sub‑structures and collective dynamical regimes.

At each scale, an analogous organisational principle is reproduced:
\begin{itemize}[label=-]
  \item certain sub‑systems realise quasi‑stationary local closures;
  \item a more flexible relational environment mediates interactions and ensures global cohesion;
  \item the hierarchical sum of the associated coherence costs defines an effective gravitational mass.
\end{itemize}

\subsection*{Example: hydrogen as the first atomic closure}
The hydrogen atom constitutes the first hierarchical manifestation in which a minimal stationary regime (the electronic state) establishes a closure on a composite structure (the proton) without destroying the internal coherence of either constituent. In this configuration, the active surface of the proton—understood as the set of surface relations that remain available for external couplings—plays a central role: it is on this logical “skin” that the electronic stationary regime stabilises.

Within this framework, the fine‑structure constant \(\alpha\) encodes the ratio between the total internal coherence of the proton and the fraction of this coherence that is projected onto its surface and remains accessible for coupling with the electron. In the classical Coulomb description, the binding potential at the Bohr radius is expressed as \(V_{\mathrm{Coulomb}}(r) = -k e^{2}/r\), which yields a ground‑state energy of approximately \(-4.36 \cdot 10^{-18}\,\mathrm{J}\). In the relational framework, one may introduce an effective potential
\[
V_{\mathrm{LDP}}(r) = -\alpha \frac{h c}{r},
\]
where \(\alpha\) quantifies this fraction of surface coherence and \(h c / r\) sets the characteristic action and propagation scale. Evaluated at the Bohr radius, this expression gives \(V_{\mathrm{LDP}}(r_{\mathrm{Bohr}}) \approx -4.35 \cdot 10^{-18}\,\mathrm{J}\). The proximity of these values in order of magnitude indicates that the relational interpretation of \(\alpha\) as a measure of surface coherence is, in principle, compatible with the observed hydrogen binding energy.

From this perspective, the conceptual boundary between “matter” and “geometry” becomes less rigid. What is usually described as matter corresponds to the portion of the coherence budget concentrated in persistent hierarchical closures; what is described as geometry corresponds to the manner in which the emergent metric is deformed under the influence of this concentration. These two descriptions thus refer to complementary aspects of a single underlying relational organisation.

\section{Indications at the cosmological scale}

On sufficiently large scales, the distribution of coherence density is no longer associated with a small number of isolated structures, but instead with the ensemble of hierarchical closures present within the observable universe. At this regime, the emergent metric becomes a genuinely global object, whose properties encode both the statistical distribution of coherence costs and the specific manner in which these costs are hierarchically organised.

A gravitational coupling constant, defined from the elementary internal structure of nucleons, sets the magnitude of the curvature induced in the global metric for a given hierarchical content. Conversely, the phenomenology observed at cosmological scales—such as the mean expansion rate, the formation of large-scale structure, and the effective curvature on cosmic distances—can be interpreted as imposing additional constraints on the spatial and temporal distribution of hierarchical closures and on the conditions under which they are sustained in cosmological proper time.

The objective here is not to construct a fully developed cosmological model, but rather to demonstrate that the relational framework admits, at least in principle, a consistent extension from the particle scale to the scale of the universe. The same conceptual entities—coherence cost, coherence density, emergent metric, and gravitational coupling—remain applicable across all levels of hierarchical complexity under consideration.

\section{Hydrogen and ordinary matter}
The hydrogen atom constitutes the first stable atomic configuration emerging from the coupling between a composite hierarchical structure (the proton) and a minimal elementary structure (the electron). Within the relational framework, it is not treated as a simple aggregate of constituent particles, but rather as the first coherent compromise between two distinct regimes of closure, thereby enabling the emergence of ordinary matter.

\subsection{Proton, electron, and the first atomic system}
In this model, the electron realises the minimal relational structure \((4,6)\), which is perfectly closed, devoid of internal hierarchy, and exhibits no coherence leakage: all pulsations remain confined within the structure. This property defines an elementary proper time that is directly encoded in its Compton frequency. By contrast, the proton is described as a composite entity with hierarchical closure, comprising a valence core, a dense relational sea, and a residual fraction of surface links, partitioned into a stationary component of the order of a few percent and a dynamic component of about \(91\%\). This configuration is associated with a substantially higher coherence cost.

The hydrogen atom results from the coupling of these two regimes within a single system: a hierarchically closed proton exhibiting a minimal leakage \(\epsilon\), surrounded by an electron of type \((4,6)\) that must preserve its own closure while remaining compatible with the active logical surface of the proton. From this perspective, the hydrogen atom is not merely the simplest atomic system; it represents the first non-trivial solution to the coexistence constraints between two stationary structures of fundamentally different natures.

\subsection{Why the electron does not fall onto the proton}

Within the framework of classical mechanics, an electron orbiting a nucleus is expected to emit electromagnetic radiation continuously, thereby losing energy and spiralling inward until it eventually collapses onto the nucleus. The very existence of a stable hydrogen atom therefore appears paradoxical unless one introduces ad hoc assumptions such as Bohr’s “non‑radiating” stationary orbits. In the LDP framework, this issue is reformulated in structural–logical terms: a collapse of the electron onto the proton would entail the destruction of the minimal \((4,6)\) closure, and thus the breakdown of the logical integrity of the corresponding elementary structure. Such a breakdown is, however, prohibited by the axioms of persistence.

In this perspective, the stability of hydrogen emerges as the result of a structural balance: the electron cannot approach the proton beyond a certain critical distance without coming into conflict with its own coherence structure, whereas the proton cannot extend its logical surface further without compromising its internal hierarchical closure. The fundamental orbit is therefore not a trajectory protected by an ad hoc dynamical postulate, but rather the distance at which both closures can be sustained simultaneously without violating their respective coherence constraints. 

\subsection{Effective potential, Bohr radius, and the role of \texorpdfstring{\(\alpha\)}{alpha}}

In the conventional formulation of atomic physics, the binding energy of the hydrogen ground state is given by \(-13.6\,\mathrm{eV}\), and the Bohr radius \(a_{0}\) sets the characteristic length scale at which the electron’s radial probability density is maximal. Within the relational framework, these quantities are reinterpreted as measures of the coherence cost necessary to sustain a \((4,6)\)-type stationary regime at a finite separation from a composite hierarchical structure, explicitly incorporating both the three‑dimensional dilution of relational links and the active logical surface associated with the proton.

The analysis demonstrates that, by reformulating the Coulomb potential on the basis of three‑dimensional closure and a finite number of surface relations \(R_{\mathrm{surf}} \simeq 502\), the fine‑structure constant \(\alpha\) naturally emerges as the ratio of surface coherence available for the electron–proton coupling. When evaluated at the Bohr radius, the resulting effective potential yields a binding energy \(E_{\mathrm{LDP}} \simeq -4.35 \times 10^{-18}\,\mathrm{J}\), which, within the adopted approximations, coincides with the standard value \(-4.36 \times 10^{-18}\,\mathrm{J}\) (i.e. \(-13.6\,\mathrm{eV}\)). This quantitative agreement supports the interpretation of \(\alpha\) as encoding a fractional surface‑coherence budget that remains consistent with empirical hydrogen spectroscopy.

\subsection{Hydrogen, complexity threshold, and ordinary matter}
The existence of hydrogen does not require postulating any additional arbitrary interaction. Rather, it relies on the possibility of stabilising a coupling between the minimal elementary structure (the electron) and the first composite hierarchical structure (the proton), in such a way that the respective closure properties of each structure are preserved. This threshold corresponds to the first form of organisation in which the internal coherence of one structure (the proton) can durably accommodate another structure (the electron) at a finite distance, by reallocating a portion of its “surface‑coherence” resources while maintaining its own structural integrity.

Ordinary matter, in its more complex manifestations, may then be interpreted as a systematic generalisation of this organisational pattern. All atomic and molecular configurations emerge from repetitions and combinations of the same basic logical scheme, in which internal hierarchical closures leave sufficiently active “surface” degrees of freedom to sustain electron‑type stationary regimes. Hydrogen thus appears as the first stable interface between two distinct levels of closure, and as the minimal necessary condition for the emergence of chemistry and, more broadly, of a structured material world on macroscopic scales.

\section{Neutron: symmetry, fatigue, and instability}
The following discussion proposes a qualitative interpretation of neutron decay within the LDP framework. It is conjectural in nature and does not, at this stage, amount to a complete quantitative derivation.

Like the proton, the neutron is a composite nucleon, but it realises a symmetric hierarchical architecture of type \((u,d,d)\), which places it at the boundary of stability: it is stable when embedded in specific nuclei, yet unstable in isolation. In the relational framework, this behavioural difference does not arise from a mere permutation of quark flavours; rather, it results from a distinct distribution of coherence cost and logical violations within the hierarchical structure.

\subsection{Relational architecture of the neutron}
Analogously to the proton, the neutron is represented by a hierarchy comprising three stationary groups (the valence quarks) and one dynamic group (the quark sea), organised according to constraints of minimal completeness and logical optimisation. The key distinction lies in the internal symmetry: whereas the proton adopts an asymmetric configuration \((u,u,d)\), the neutron realises a symmetric structure \((u,d,d)\), which modifies both the distribution of violations and the compensatory capacity of the quark sea.

In this description, the groups associated with the \(u\) quark and the two \(d\) quarks contain numbers of entities and relations comparable to those attributed to the proton (for example, 24 relations for the \(u\) group and 28 for each \(d\) group, with a quark sea comprising on the order of \(10^{4}\) dynamic relations). However, the duplication of the \(d\) group imposes a more rigid internal symmetry. The total number of internal relations thus reaches \(10\,992\), with a stationary fraction of approximately \(9\%\) and a dynamic fraction of approximately \(91\%\), as in the proton, but distributed such that the \(d\) quarks occupy a more constrained role in achieving global closure.

\subsection{Compton wavelength and logical compression}
As in the case of the proton, the global structure of the neutron can be analysed in terms of its Compton wavelength, which encodes the characteristic frequency associated with its rest mass. Using the experimental value \(\lambda_{C,n} \simeq 1.31959 \times 10^{-15}\,\mathrm{m}\), one can define a “logical” radius associated with the stationary, closed-loop wave and compare it with an effective “matter” radius \(R_{n} \simeq 0.87 \times 10^{-15}\,\mathrm{m}\), taken from nuclear data.

The calculation shows that the radius inferred from the neutron Compton wavelength is smaller than the adopted matter radius, with a relative deviation of approximately \(3.56\%\). This suggests that the closure wave induces a non-negligible compression of the corresponding hierarchical structure. Performing the same comparison for the proton yields a relative deviation of about \(0.07\%\), i.e. more than fifty times smaller than for the neutron, which indicates that the neutron structure is subject to a much stronger rhythmic constraint.

\subsection{Closure fatigue and neutron decay}
Within the LDP framework, this logical compression can be interpreted as an increased coherence cost required to maintain the hierarchical closure of the neutron at each cycle of global pulsation. The internal symmetry \((u,d,d)\) distributes violations less favourably than the asymmetric configuration of the proton, such that the quark sea must compensate a larger fraction of the internal tensions, thereby increasing the overall rigidity of the structure.

Over long proper times, this situation can be viewed as a process of closure fatigue: although residual violations are minimised at each cycle, they cannot be completely eliminated without a hierarchical reconfiguration, which gradually reduces the ability of the form to maintain its coherence. The decay of a free neutron \((n \rightarrow p + e^{-} + \bar{\nu}_{e})\) thus appears as the resolution of this impasse: the symmetric structure yields and reorganises into a more stable proton, accompanied by an electron and an antineutrino that carry away the fraction of coherence that has become incompatible with the initial closure.

\subsection{Neutron, nuclei, and ordinary matter}
The unstable nature of the free neutron contrasts with its effective stability when embedded in many nuclei, where it can persist indefinitely provided that the global nuclear configuration remains compatible with the relevant coherence constraints. Within the relational framework, this stabilisation is interpreted as a collective phenomenon: the presence of additional nucleons redistributes local violations and allows the intrinsic compression of the neutron to be partially offloaded onto the emergent metric and the surrounding closures, thereby mitigating internal fatigue.

The behavioural difference between proton and neutron is therefore not a secondary compositional detail, but rather the manifestation of two distinct organisational schemes for composite hierarchical closures constructed under the same minimal axioms. At macroscopic scales, this asymmetry provides a qualitative explanation for the predominance of protons in observable ordinary matter (in association with electrons to form atoms), while neutrons primarily assume the role of structural binders within nuclei, subject to a more delicate trade-off between local stability and global fatigue.

The mechanism proposed here offers a qualitative account of neutron instability within this relational perspective; however, the corresponding estimates of the neutron lifetime remain highly approximate and do not supersede standard calculations based on the weak interaction.

\section{Photon: ephemeral triangular configuration and information transport}

The identification of the photon with a minimal \((3,3)\)-type structure must be regarded as a highly speculative working hypothesis. It is not, at this stage, deduced from the axiomatic basis of the model and does not, in its current formulation, recover the full set of empirically established properties of the photon (such as polarisation states, spin, or the quantum electrodynamics formalism).

Within the LDP framework, the photon occupies a singular status: it is neither a closed volumetric entity, as in the case of the electron, nor a hierarchical composite structure, as for nucleons. Rather, it is construed as the minimal surface motif capable of transmitting information and energy between two closures. Its evanescent character follows directly from the axioms: a purely logical triangle is unable to stabilise as a volume; it can only exist for the duration of a transition between two coherent configurations.

\subsection{Minimal logical triangle}

In the relational architecture under consideration, three mutually connected entities constitute the first non-trivial surface motif: a triangle. This motif can sustain a form of “tension” but cannot close into a stable three-dimensional volume. While four entities (a tetrahedron \((4,6)\)) realise the minimal volumetric closure, thereby giving rise to a persistent structure (identified here with the electron), three entities generate only a transient consistency, intrinsically exposed to dissipation.

The LDP hypothesis therefore proposes to associate the photon with this temporary agreement among three entities oscillating in a perfectly triangular configuration, “an instant of pure surface, without shoulders to bear the weight of reality over time”. This triangle possesses neither internal hierarchy nor an internal dynamic medium, and it does not constitute a complete closure. It is entirely sustained by the triangular tension mandated by axiom 2 and is, for that reason, intrinsically unstable once it is no longer maintained by the transition process that engendered it.

\subsection{\texorpdfstring{\(E = \hbar \omega\)}{E = ħω} re-read in LDP}

In conventional physics, the relation \(E = \hbar \omega\) assigns to each photon an energy proportional to its frequency, typically without providing a more detailed structural interpretation. Within the LDP framework, this relation is reinterpreted as the quantitative expression of the coherence cost necessary to construct and sustain, for a finite duration, the logical triangle that mediates the transition between two closures.

Each photon is associated with a transition between two stationary regimes (for instance, two bound electronic levels or two nuclear states), and the frequency \(\omega\) encodes the temporal rhythm of this transition within the emergent metric. The reduced Planck constant \(\hbar\), previously interpreted as the minimal granularity of action characterizing a closed stationary regime, now quantifies the amount of action per cycle that is made available to the triangular surface; the photon energy thereby corresponds to the product of this action granularity by the cadence imposed on the ephemeral logical motif.

\subsection{Compton effect: coherence test}

The Compton effect constitutes a direct testing ground for this relational interpretation of the photon. In the standard description, a photon of incident wavelength \(\lambda_{\mathrm{inc}}\) scatters off an almost free electron and emerges with an increased wavelength \(\lambda_{\mathrm{diff}}\), according to
\[
\lambda_{\mathrm{diff}} - \lambda_{\mathrm{inc}} = \lambda_{C,e} (1 - \cos\theta),
\]
where \(\lambda_{C,e}\) denotes the Compton wavelength of the electron and \(\theta\) the scattering angle.

In the LDP interpretation, the electron is modeled as a perfectly coherent closed tetrahedron \((4,6)\) (\(\epsilon = 0\)), supporting a stationary wave at the Compton frequency, while the photon is represented as a rhythmic triangle that tangentially interacts with the logical surface of this electron. The photon impact is described as a local modulation of the electron’s stationary wave:
\begin{itemize}[label=-]
    \item the triangular violation \(\nu\) temporarily increases on a localized portion of the surface (local transition from \(\nu = 0\) to \(\nu > 0\));
    \item a fraction of the photon’s rhythm is transferred to the internal pulsation of the electron, which registers a logical impulse (the “recoil”);
    \item the triangle subsequently emerges with a reduced frequency \(\omega' < \omega\), corresponding to the increased wavelength, while the global closure of the electron is reestablished (\(\nu \rightarrow 0\) after relaxation).
\end{itemize}

The conserved quantities of the standard treatment (energy and momentum) are then reinterpreted as the conservation of the action budget \(\hbar\) and of surface coherence on the electronic sphere: the variation in photon frequency is exactly equal to that required to preserve the integrity of the \((4,6)\) structure while accommodating the induced perturbation.

\subsection{Electromagnetic spectrum and transitions}
Within this framework, the electromagnetic spectrum is not conceived as an arbitrary collection of admissible frequencies, but as the set of transitions that are compatible with the coherence constraints imposed by the closures involved. Each observable photon is associated with a transition in which two stationary structures (atomic, nuclear, or of higher complexity) undergo a reconfiguration that preserves their respective internal closures, while externalising the resulting coherence difference in the form of a logical triangle characterised by a frequency \(\omega\).

The relation \(E = \hbar \omega\) ensures that any energy difference between two stationary states can be represented by a single triangular surface, without the need to introduce additional parameters: the photon frequency is fully determined by the nature of the transition, and the adaptability of the triangular motif makes it possible, in principle, to account for all empirically observed frequencies, from radio waves to gamma rays. In this perspective, the photon appears as the minimal, mobile unit of the electromagnetic spectrum, not as a particle in the usual volumetric sense, but as an evanescent contract established between two closures, which vanishes as soon as the logical tension it conveys is absorbed or redistributed within the relational network.

This construction must therefore be regarded as a preliminary and exploratory proposal rather than as a complete theory: a fully developed formulation should demonstrate how Maxwell’s equations and the formal structure of quantum electrodynamics (QED) can be derived from, and rigorously embedded within, the underlying relational framework.

\section{Constants, measurable effects, and status of the LDP}
Within the LDP framework, physical constants are not treated as external parameters adjusted a posteriori to fit experimental data. Instead, they are regarded as quantitative indicators of distinct forms of coherence granularity: action granularity for an internal pulsation cycle, surface-coupling granularity for composite structures, thermal granularity for statistical ensembles, and gravitational leakage granularity toward the global relational background. This section compiles the principal constants employed in the text and clarifies their conceptual status within the relational description. 

\subsection{\texorpdfstring{\(\hbar\)}{ħ}: minimal action granularity}
In the relational formulation, the reduced Planck constant \(\hbar\) is interpreted as the minimal quantum of action associated with an elementary logical stationary regime. Each internal pulsation cycle of a minimal \((4,6)\)-type structure, identified here with the electron, carries a discrete unit of action of magnitude \(\hbar\). When a physical realisation of such a pulsation is available, the proper energy of the structure is expressed as \(E = \hbar \omega\). 

Under this interpretation, the relation \(E = \hbar \omega\) does not represent the addition of energy to a pre-existing substrate, but rather the coherence cost required to sustain an elementary stationary regime. The constant \(\hbar\) thereby determines the minimal increment by which internal coherence can self-close. It is thus understood as a constant of logical action granularity, rather than as an empirically tuned, external parameter.

\subsection{\texorpdfstring{\(\alpha\)}{α} and \texorpdfstring{\(\varphi\)}{φ}: surface coherence and geometric optimisation}
Within this framework, the fine-structure constant \(\alpha\) is interpreted in an analogous manner to the preceding quantities: not as a parameter without origin, but as a quantitative measure of surface-coupling granularity. Specifically, it characterises the fraction of the internal coherence budget of a proton that can effectively be expressed at its surface when a stable coupling with an electron is established. 

A first-order estimate, based on an effective active surface of magnitude \(R_{\mathrm{surf}} \simeq 500\) relations, already produced a value close to \(\alpha^{-1}\) via the relation \(\alpha^{-1} \simeq R_{\mathrm{surf}}^{2} / \mu\), where \(\mu = m_{p}/m_{e}\) denotes the proton–electron mass ratio. A refinement of this estimate reveals a more structured picture: the total number of stationary valence relations in the proton is taken to be \(r_{\mathrm{valence}} = 930\), and one third of this quantity, \(r_{\mathrm{valence}}/3 = 310\), is interpreted as the average coherence core per valence quark. 

Multiplying this third by the golden ratio \(\varphi = (1 + \sqrt{5})/2\) yields a refined estimate of the active surface
\[
R_{\mathrm{surf}}^{(\varphi)} \simeq \varphi \frac{r_{\mathrm{valence}}}{3} \approx 501.59,
\]
which is constructed entirely from parameters internal to the protonic architecture together with a simple geometric factor. Substituting this expression into \(\alpha^{-1} \simeq \left(R_{\mathrm{surf}}^{(\varphi)}\right)^{2}/\mu\), with \(\mu \simeq 1836.15\), leads to \(\alpha^{-1} \simeq 137.022\), in agreement with the experimental value \(\alpha^{-1} \simeq 137.036\) to within a few \(10^{-4}\) in relative terms, without the introduction of any ad hoc correction. 

Accordingly, the fine-structure constant can be reformulated as
\[
\alpha^{-1} \simeq \frac{\varphi^{2}}{\mu} \left(\frac{r_{\mathrm{valence}}}{3}\right)^{2},
\]
an expression written solely in terms of \(\mu\), \(r_{\mathrm{valence}}\), and the geometric factor \(\varphi\). In this perspective, the golden ratio is not invoked as a merely aesthetic motif, but proposed as a candidate constant of relational geometric optimisation: it describes the most efficient redistribution of the coherence content of an average stationary core toward the active surface, subject to constraints of hierarchical closure and minimal coherence leakage. 

This interpretation is compatible with an intuition already present in the literature, according to which the fine-structure constant relates a dimensionless mass ratio \(\mu\) to a finite number of coupling channels between elementary entities. In the present framework, this idea is reformulated in terms of active surfaces and relational topology: \(\alpha\) quantifies the fraction of the action granularity set by \(\hbar\) that remains available at the surface for interaction channels, once internal hierarchical closure has been ensured.

\subsection{\texorpdfstring{\(k_{B}\)}{kB}: thermal granularity and protonic hierarchy}
The Boltzmann constant \(k_{B}\) assumes an analogous role once one considers, instead of an isolated structure, a statistical ensemble of hierarchical structures. In a relational framework, the mean thermal energy \(E_{\mathrm{th}} \sim k_{B} T\) quantifies the fraction of the global coherence budget that is, on average, mobilised in the collective excitations spanning a large number of pulsation cycles. 

In this perspective, temperature is no longer treated as a primitive variable: it characterises the effective occupation of excited coherence cycles within an ensemble of structures (electrons, nucleons, atoms). The function of \(k_{B}\) is to map this statistical population counting to a corresponding average energy scale, in a manner analogous to the way in which \(\hbar\) connects an elementary cycle to a quantum of action. Accordingly, \(k_{B}\) sets the characteristic energy scale of a thermal fluctuation per degree of freedom. 

In the current formulation of the LDP, \(k_{B}\) is not yet uniquely derived, but a candidate expression is inferred from the combined architectures of the electron and the proton. This expression couples: the minimal rhythmic tension \(h/\tau_{0}\) associated with a single pulsation cycle of the electron; the number of active relations of the \((4,6)\) structure; the electron–proton mass ratio \(\mu\); the dynamical fraction \(R_{\mathrm{mer}}/R_{\mathrm{tot}} \simeq 91\%\) of the quark sea; and the leakage rate \(\varepsilon\), which quantifies the fraction of coherence dissipated into the global relational background. 

Under these hypotheses, one obtains an expression that introduces no additional continuous free parameter,
\[
k_{B}^{\mathrm{LDP}} = \frac{h}{\tau_{0}} \, \frac{n_{\mathrm{rel}}^{e}}{\mu^{2}} \, \frac{R_{\mathrm{mer}}}{R_{\mathrm{tot}}} \, \frac{\varepsilon}{4},
\]
where each factor is associated with a specific property of the electronic and protonic architectures. Numerically, this relation reproduces the empirical value of \(k_{B}\) with a controlled relative deviation, suggesting that the same relational hierarchy that fixes \(\alpha\) may also govern the emergent thermal energy scale.

\subsection{\texorpdfstring{\(G\)}{G} and \texorpdfstring{\(\varepsilon\)}{ε}: effective gravitational coupling}

Within the LDP framework, the Newtonian gravitational constant \(G\) is understood as the metric manifestation of a more elementary coupling mechanism: the minimal coherence leakage that connects hierarchical structures to the global relational background. The internal hierarchy of a proton cannot simultaneously realise the maximal closure of all its subsystems and the global closure of the composite system as a whole. This incompatibility induces a non-vanishing but extremely small coherence leakage \(\varepsilon\) toward the surrounding relational framework.

The parameter \(\varepsilon\) is defined as the relative deviation between an ideal hierarchical closure and the effective closure actually realised by the protonic structure, given its \(11\,017\) internal relations and the partition between stationary and dynamical components. It thereby establishes a quantitative correspondence between (i) the internal relational architecture of the proton, (ii) the curvature of the emergent metric, and (iii) an effective gravitational coupling constant. The latter may be interpreted as a derived quantity, \(G_{\mathrm{LDP}}\), rather than as a fundamental constant introduced \emph{a priori}.

At the present stage, the explicit form of \(G_{\mathrm{LDP}}\) remains conjectural and is only partially constrained through comparison with cosmological scales (for instance via the Hubble parameter). Its conceptual status is nonetheless unambiguous: it is not a freely adjustable parameter introduced to reproduce an observed interaction strength,

\section{Perspectives and possible tests}
The LDP framework is not conceived as a direct competitor to established theories within their respective domains of validity; rather, it proposes a reinterpretation of their fundamental entities (particles, constants, metric) as emergent manifestations of a minimal logic of coherence. From this standpoint, the primary issue is not whether the framework immediately reproduces all known numerical values with high experimental accuracy, but whether it gives rise to a testable research programme capable of falsifying or refining its underlying combinatorial hypotheses.

The preceding sections have presented several quantitative indications:
\begin{itemize}[label=-]
    \item the association of the \((4,6)\) structure with the electron and of the (\(24\),\(28\),\(930\),\(10\,087\),\(11\,017\)) hierarchy with the proton;
    \item the reinterpretation of the fine-structure constant \(\alpha\) in terms of an active surface \(R_{\mathrm{surf}}^{(\varphi)} \simeq \varphi \, r_{\mathrm{valence}}/3\), yielding agreement at the level of better than a few \(10^{-4}\) in relative value;
    \item a candidate expression for \(k_{B}\), constructed from the same integers, the dynamical fraction of the sea, and the leakage rate \(\varepsilon\), which is compatible with the measured value within a controlled error budget;
    \item an interpretation of the effective gravitational coupling in terms of minimal coherence leakage, calibrated to reproduce cosmological orders of magnitude.
\end{itemize}
These numerical agreements should not be interpreted as precise predictions in the strict sense, but as indicators of internal coherence: they suggest that a common set of combinatorial building blocks may underlie several dimensionless constants. The residual discrepancies, typically in the range \(10^{-4}\) to \(10^{-2}\), are not to be regarded as insignificant corrections; on the contrary, they delineate a potential testing strategy. If these deviations prove to be systematic and mutually correlated, they may encode coherence-leakage or hierarchical-structure effects that are not yet modelled with sufficient resolution.

\subsection{Tests on the \texorpdfstring{\((4,6)\)}{(4,6)} structure and internal hierarchies}

A first line of investigation concerns the robustness of the minimal \((4,6)\) structure and of the associated protonic hierarchy \((24,28,930,10087,11017)\). Within the LDP framework, these integers are not introduced ad hoc; rather, they arise from constraints of minimal completeness, triangular coherence, and logical optimisation. Any empirical refutation of this structural architecture would therefore directly undermine the validity of the entire framework.

Several concrete research directions can be specified:
\begin{itemize}[label=-]
    \item \textbf{High-precision spectroscopy of simple bound systems} (hydrogen, muonic hydrogen, hydrogen-like ions): if the fine-structure constant \(\alpha\) is indeed functionally related to \(r_{\mathrm{valence}}\), \(\varphi\), and \(\mu\), then fine and hyperfine structure corrections should be reinterpretable as deformations of the surface architecture \(R_{\mathrm{surf}}^{(\varphi)}\) rather than as independent phenomenological parameters. Small deviations from standard theoretical predictions could then be used to extract effective estimates of higher-order hierarchical corrections that are not yet incorporated into the current formulation.
    \item \textbf{Extensions to heavier nucleons and nuclei}: the present analysis has concentrated on the proton as the minimal composite configuration. A natural extension consists in applying the same combinatorial procedure to other nucleons (e.g. the neutron) and to light nuclei, in order to identify hierarchical architectures that satisfy the LDP axioms while simultaneously reproducing observed mass ratios, charge radii, and stationary/dynamical energy fractions. Should the integers selected for these extended structures fail to preserve the central role of \(\varphi\) and \(\varepsilon\), the proposed construction of \(\alpha\), \(k_{B}\), and \(G\) would need to be critically re-evaluated.
    \item \textbf{Combinatorial simulations}: systematic numerical exploration of the configuration space \((n_{u}, n_{d}, r_{\mathrm{valence}}, R_{\mathrm{mer}})\), under the constraints of coherence and minimality, can be employed to test whether the pair \((24,28)\) and the total \(r_{\mathrm{valence}} = 930\) are indeed singled out as optimal, and whether alternative architectures fail to reproduce, in a unified manner, the electric charge, the dynamical fraction, and the empirical mass ratios.
\end{itemize}

\subsection{\texorpdfstring{\(\alpha\), \(\varphi\)}{α, φ} and residual discrepancies}
The reinterpretation of \(\alpha^{-1}\) as
\[
\alpha^{-1} \simeq \frac{\varphi^{2}}{\mu} \left(\frac{r_{\mathrm{valence}}}{3}\right)^{2}
\]
is based on a specific definition of the active surface \(R_{\mathrm{surf}}^{(\varphi)}\) and on the hypothesis that \(\varphi\) functions as a relational geometric optimisation parameter. The resulting numerical agreements are promising; however, the presence of a residual discrepancy, however small, is of central importance: within the LDP framework, it is not merely acceptable but is expected to carry physical meaning.

A natural approach is to interpret these discrepancies as signatures of surface corrections induced by the leakage parameter \(\varepsilon\) and by the dynamics of the relational sea. In the present construction, \(R_{\mathrm{surf}}^{(\varphi)}\) is expressed using exact integers \((930)\) and \(\varphi\), and then compared with the constant \(\alpha\) as determined in specific experimental regimes (bound states, quantum vacuum, renormalisation). It is therefore plausible that the residual deviation:
\begin{itemize}[label=-]
    \item encodes contributions from hierarchical levels not yet incorporated into the model (for instance, collective modes on larger spatial or relational scales);
    \item or reflects the fact that the effective active surface is “rough” or fluctuating, rather than ideally smoothed by the action of \(\varphi\).
\end{itemize}
A corresponding research programme would involve:
\begin{itemize}[label=-]
    \item systematically comparing the value \(\alpha_{\mathrm{LDP}}\), inferred from \(R_{\mathrm{surf}}^{(\varphi)}\), with independent experimental determinations of \(\alpha\) as a function of the energy scale, in order to assess whether the residuals do or do not track the same behaviour as the renormalisation-group “running” of \(\alpha\);
    \item theoretically analysing how surface corrections associated with \(\varepsilon\) or with the dynamical fraction of the relational sea could perturb \(R_{\mathrm{surf}}^{(\varphi)}\) by an amount quantitatively consistent with the observed discrepancies.
\end{itemize}
In this perspective, the most stringent tests of the LDP are not provided by the agreement of the central values alone, but by the detailed structure of the residuals: their magnitude and their correlations with other derived constants.

\subsection{\texorpdfstring{\(k_{B}\)}{kB}: thermal scale and relational statistics}
The candidate expression for \(k_{B}\) proposed in this work connects:
\begin{itemize}[label=-]
    \item the electron’s intrinsic pulsation (via \(h/\tau_{0}\));
    \item the number of relations characterizing the minimal \((4,6)\) structure;
    \item the mass ratio \(\mu\);
    \item the dynamical fraction \(R_{\mathrm{mer}}/R_{\mathrm{tot}}\) associated with the proton;
    \item and the leakage rate \(\varepsilon\), interpreted as the fraction of available coherence that can participate in collective excitations.
\end{itemize}
This construction yields a thermal scale consistent with the experimentally determined value of \(k_{B}\), without the introduction of any additional continuous parameter. As before, the remaining discrepancies may be exploited as a quantitative probe for testing the framework.

Several research directions follow from this proposal:
\begin{itemize}[label=-]
    \item \textbf{Thermodynamics of simple systems}: re-examine the heat capacities of model systems (such as a weakly interacting electron gas or a nucleon gas in a regime where nuclear interactions remain under quantitative control) in terms of the occupation of coherence cycles, and assess whether the expression \(k_{B}^{\mathrm{LDP}}\) provides an accurate description of the crossover between quantum and classical regimes.
    \item \textbf{Role of \(\varepsilon\)}: since \(k_{B}^{\mathrm{LDP}}\) depends explicitly on \(\varepsilon\), any refinement in the estimation of the leakage rate (for instance, derived from gravitational or cosmological analyses) should induce a corresponding correction in the inferred value of \(k_{B}\). A systematic mismatch between these two determinations of \(\varepsilon\) would furnish a direct empirical test of the proposed framework.
    \item \textbf{Fluctuations and entropy}: the standard relation \(S = k_{B} \ln W\) can be reformulated in terms of the number of relational configurations accessible to a given ensemble of hierarchical structures. A more detailed investigation, based on combinatorial models, could determine whether the scale \(k_{B}^{\mathrm{LDP}}\) is indeed the appropriate one to connect the count of relational configurations with macroscopic thermodynamic quantities.
\end{itemize}

\subsection{Gravitation, leakage \texorpdfstring{\(\varepsilon\)}{ε}, and cosmological tests}
On the gravitational side, the LDP proposes to reinterpret the curvature of the emergent metric as the macroscopic trace, within an effective description, of the underlying distribution of relational density and of the associated coherence cost. The minimal leakage \(\varepsilon\), defined from the closure constraints of the proton, plays the role of a coupling parameter between local hierarchical structures and the global relational background.

Several classes of potential tests can be envisaged:
\begin{itemize}[label=-]
    \item \textbf{Consistency with general relativity}: in regimes where general relativity has been stringently tested (solar-system observations, binary pulsar timing, propagation of gravitational waves), the objective is to verify that the relational reinterpretation provides at least a conceptually coherent “lift” of the standard framework, without altering the numerical predictions. Any systematic discrepancy would indicate the need to refine the mapping between \(\varepsilon\) and the effective curvature.
    \item \textbf{Cosmological orders of magnitude}: the same parameter \(\varepsilon\) also appears in attempts to relate protonic microstructure to cosmological scales (for example via the effective energy density or the Hubble parameter). If such connections are empirically supported, adjusting a single parameter \(\varepsilon\) could, in principle, account simultaneously for the strength of local gravitational coupling and for selected cosmological orders of magnitude.
    \item \textbf{Residual deviations in \(G\)}: the gravitational constant \(G\) is known with significantly lower precision than other fundamental constants such as \(\alpha\) or \(k_{B}\). Within the LDP framework, experimental dispersions in the measured values of \(G\) are interpreted not solely as metrological limitations, but also as a possible probe of fine relational structure: apparent variations could encode differences in effective relational density or in the hierarchical composition of the probed systems.
\end{itemize}

\subsection{Status of the LDP with respect to observations}
The LDP currently remains at the stage of a foundational theoretical framework. Several of its constructions are still conjectural; certain expressions (in particular those proposed for \(k_{B}\) and \(G\)) should be regarded as structured working hypotheses rather than definitive derivations; and its treatment of out-of-equilibrium phenomena and of detailed cosmological dynamics is, at present, only partial. Nevertheless, the framework yields a set of convergent empirical tests:
\begin{itemize}[label=-]
    \item verify that the same relational integers and the same factor \(\varphi\) can organise \(\alpha\), \(k_{B}\), and \(G\) across distinct observational regimes;
    \item interpret residual discrepancies as independent estimates of the leakage parameter \(\varepsilon\) and of hierarchical corrections, rather than as arbitrary or purely systematic errors;
    \item extend the construction to structures beyond the electron and the proton, while preserving the internal coherence of the role attributed to \(\varphi\) as a relational geometric optimisation invariant.
\end{itemize}
From this standpoint, the scientific value of the LDP is not assessed solely on the basis of the numerical proximity of its predictions to the measured values of the constants, but also on the manner in which each discrepancy, however small, is converted into a precise question regarding the underlying structure of reality: which component of the hierarchy, of the active surface, or of the coherence leakage has not yet been incorporated into the description?

\section{Conjectural genesis of closures}
The LDP framework is deliberately formulated in an a-cosmological manner: it specifies the logical conditions for the existence of closed structures without presupposing any particular cosmological history. Within this perspective, stable forms (electron, proton, atomic nuclei, etc.) are construed as solutions to optimisation problems under constraints, which can be instantiated wherever a sufficiently structured relational background is present. It remains, however, to outline a scenario in which these closures do not appear as given from the outset but instead arise as the outcome of large-scale constraint events. 

One may conjecture that, underlying any metric reality, there exists a pre-physical logical field, analogous to a matrix of pulsating informational elements, in which only axioms of minimal pulsation and coherence are in force. In this field, elementary pulsations do not yet instantiate any hierarchical organisation; they merely realise the logical possibility that “something occurs rather than nothing”. As long as the relational density remains low, configurations are essentially transient and do not generate durable closures. Only when regions of the field are subjected to sufficiently intense collective constraints do more articulated compositions emerge: the majority remain unstable, but some attain minimal closure (photons, quarks, then electrons and protons) and persist. 

From this standpoint, what metric cosmology characterises as the Big Bang can be interpreted as an episode during which a finite region of the pre-physical logical field crosses a constraint threshold, thereby inducing the rapid selection of closed architectures capable of sustaining a stable regime of reality. It is not necessary to posit the uniqueness of such an episode: nothing in the LDP prohibits multiple regions of the relational background from undergoing, at different “moments” in the emergent sense, analogous transitions toward families of closed forms. From our internal observational standpoint, we would have access only to the cosmological branch to which our own closures belong.

Black holes may thus be reconsidered, still at a conjectural level, as extreme constraint configurations in which stabilised relational hierarchies can no longer be sustained in the form of composite structures. In the vicinity of a critical relational density, the internal hierarchy of forms disintegrates: global closure becomes incompatible with triangular coherence, and the available “coherence budget” is released in the form of once again weakly structured, pulsating elements. Within such a framework, what appears, from the standpoint of our metric description, as the “disappearance” of matter into a black hole could correspond, in LDP terms, to a recycling of coherence into the pre-physical field, potentially capable of initiating another Big-Bang-type episode elsewhere. 

A more radical hypothesis is to conceptualise this global process as forming a loop without a beginning or an end in the strict sense. In the LDP framework, time exists only as the metric translation of the counting of coherence cycles in closed structures. Outside these structures, within the pre-physical logical field, there would be neither “before” nor “after”, but solely constraint and coherence conditions. From this perspective, references to a “before the Big Bang” or an “after the evaporation of a black hole” are meaningful only with respect to stationary regimes that already sustain an emergent notion of proper time; the relational background itself would remain atemporal. 

This genesis-and-recycling scenario for closures is explicitly speculative. It does not purport to replace established cosmological models, nor to yield a quantitative reconstruction of the history of the observable universe. Its function is more modest yet complementary: to indicate that the LDP framework naturally accommodates a view in which forms do not achieve closure “for free”, but rather under the action of constraint events within a more primitive pulsating field, and in which extreme gravitational singularities could contribute to the recycling of coherence. This cosmological outline is not presented as a prediction, but as an invitation to investigate how LDP tools might eventually be integrated with more detailed models of the origin and ultimate fate of physical structures. 

The present section therefore does not have the status of a quantitative cosmological model; instead, it is proposed as a possible reinterpretation, in LDP terms, of scenarios already discussed in cosmology (such as cyclic universes or universes emerging from black holes). Should this picture be corroborated, the same relational architecture that structures \(\hbar\), \(\alpha\), \(k_{B}\), and \(G\) could also impose constraints on the manner in which Big-Bang-type episodes and extreme gravitational singularities are encoded in, and translated back into, the relational background.

\newpage

\section*{Conclusion}
\addcontentsline{toc}{section}{Conclusion}
The Projective Dynamic Logo (LDP) was constructed on a conceptually simple yet technically demanding hypothesis: to reconstruct the very possibility of physical reality from purely relational axioms, without presupposing the prior existence of space, time, or particles. Within this framework, existence is no longer grounded in an underlying substance or support, but in the capacity of certain structures to instantiate and sustain a stable distinction within a field of elementary pulsations. The minimal \((4,6)\) structure, deduced from the axioms of pulsation, triangular coherence, completeness, and optimization, thereby assumes the role of an elementary building block: it realizes the first admissible logical closure and furnishes a prototype for a relational interpretation of the electron.

On this basis, the text has shown how more complex organizations can emerge through the superposition and hierarchical arrangement of these minimal closures. In particular, the proton has been modeled as a minimal composite architecture that combines three local stationary regimes (valence quarks), a dynamically fluctuating relational sea, a finite active surface, and a coherence leakage \(\varepsilon\) into the global relational background. The integers characterizing this architecture, \(24, 28, 930, 10087, 11017\), are not introduced as free phenomenological parameters: they are selected by the logical optimization functional under stringent combinatorial constraints (multiples of 4, locally complete graphs, dominant dynamical fraction, net charge \(+1\)). They thereby provide a concrete foundation for reinterpreting several fundamental constants as signatures of internal coherence, rather than as quantities externally imposed on the structure.

Within this relational perspective, the reduced Planck constant \(\hbar\) no longer appears as an empirical parameter without theoretical origin: it quantifies the minimal granularity of action associated with a single pulsation cycle of an elementary stationary regime of type \((4,6)\), and the relation \(E = \hbar \omega\) is interpreted as the coherence cost of sustaining such a regime. The fine-structure constant \(\alpha\) admits an analogous reading, but at the scale of the proton surface: it measures the fraction of the internal coherence budget that remains available at the surface for external couplings. The refinement introduced by the golden ratio \(\varphi\), via the active surface
\[
R_{\operatorname{surf}}^{(\varphi)} \simeq \varphi \,\frac{r_{\mathrm{valence}}}{3},
\]
leads to the expression
\[
\alpha^{-1} \simeq \frac{\varphi^{2}}{\mu} \left(\frac{r_{\mathrm{valence}}}{3}\right)^{2},
\]
which reproduces the experimental value of \(\alpha\) to better than a few \(10^{-4}\) in relative deviation, without the introduction of ad hoc fitting parameters. In this context, the golden ratio emerges as a candidate invariant of relational geometric optimization: it functions as a factor encoding the optimal distribution of coherence between the stationary core and the active surface, rather than as a numerological pattern imposed from outside the theory.

The same logical architecture has made it possible to formulate candidate expressions for the Boltzmann constant \(k_{B}\) and for an effective gravitational coupling \(G_{\mathrm{LDP}}\). In the relational interpretation, \(k_{B}\) quantifies a thermal granularity, understood as the characteristic energy scale associated with the statistical excitation of multiple coherence cycles constructed from the electron’s intrinsic pulsation, the mass ratio \(\mu\), the dynamical fraction of the quark sea, and the leakage rate \(\varepsilon\). In turn, the gravitational coupling is construed as the metric transcription of the minimal coherence leakage required by the impossibility, for a composite hierarchical structure, of simultaneously saturating all local and global closure conditions. In all these instances, the numerical correspondences obtained are only approximate; however, they rely on the same combinatorial ingredients: the \((4,6)\) structure, the protonic architecture, the ratio \(\mu\), the factor \(\varphi\), and the parameter \(\varepsilon\). 

The present text does not, however, claim to provide a complete theory of fundamental constants, gravitation, or thermodynamic phenomena. Several constructions—particularly those concerning \(k_{B}\), \(G_{\mathrm{LDP}}\), and cosmology—remain conjectural or are partially constrained by empirical data, and the exploration of the space of combinatorial configurations has so far been limited. The scientific contribution of the LDP should therefore be evaluated less in terms of numerical precision than in terms of the internal coherence of the framework it proposes: a theoretical setting in which dimensionless constants, particle structures, emergent metric properties, and proper time are interpreted as distinct manifestations of a single minimal logic of coherence. Residual discrepancies relative to experimental measurements, which are far from negligible details, are instead treated as a parameter space within which to test and refine the roles of the active surface, the internal hierarchy, and the leakage parameter \(\varepsilon\). 

Finally, the section devoted to the conjectural genesis of closures outlines a possible cosmological integration of the LDP. In this scenario, a pre-physical logical field of elementary pulsations furnishes an atemporal substrate from which, under global coherence constraints, stable closed structures emerge locally, potentially via Big-Bang-like episodes. Black holes are reinterpreted as extreme-constraint configurations in which stabilized hierarchies break down, thereby recycling coherence into the relational background rather than producing an opaque geometric singularity. This cosmological perspective is explicitly presented as speculative: it is not proposed as a fully quantitative model in competition with established cosmological frameworks, but rather as a possible reinterpretation of their main qualitative features in the conceptual language of the LDP, consistent with the relational treatment of fundamental constants.

We have now established a formal translation of the LDP axioms into a framework based on signed graphs and discrete dynamical systems. Within this formalism, the \((4,6)\) configuration emerges as the first elementary stationary structure, while the \(u\), \(d\) blocks and the protonic structure appear as hierarchical closures selected by the same principles of triangular coherence, completeness, and optimization. Nonetheless, the detailed combinatorial analysis of internal blocks—particularly the proof of the exact minimality of the integers \(24, 28, 276,\) and \(378\)—as well as the rigorous derivation of certain dimensionless ratios, remain open problems. Addressing these questions will require more systematic mathematical and numerical investigations. These limitations should be interpreted as opportunities to deepen and empirically test the framework, rather than as definitive shortcomings.

The Projective Dynamic Logo should therefore be regarded as a foundational framework in the process of formalization. It provides a precise method for formulating the question of existence, for deriving a minimally coherent structure, and, on this basis, for articulating the concepts of particle, dimensionless constant, proper time, effective gravitation, and cosmological genesis. Future work will focus on refining the constructions presented here, exploring the space of hierarchical architectures in a more systematic manner, and, above all, increasing the number and precision of quantitative points of contact with experiment, so that the LDP may be corroborated, revised, or invalidated by the very phenomena it seeks to reinterpret.

\end{document}